\documentclass[12pt,a4paper]{article}
\usepackage[spanish,es-tabla,es-noquoting]{babel}
\usepackage[utf8]{inputenc}
\usepackage[T1]{fontenc}
\usepackage{amsfonts}

\renewcommand{\rmdefault}{ptm}
\usepackage{mathptmx}
\DeclareSymbolFont{Xlargesymbols}{OMX}{cmex}{m}{n}
\DeclareMathSymbol{\Xsum}{\mathop}{Xlargesymbols}{80}

\DeclareSymbolFont{newtxletters}{OML}{ntxmi}{m}{it}
\SetSymbolFont{newtxletters}{bold}{OML}{ntxmi}{b}{it}
\DeclareMathSymbol{\mu}{\mathord}{newtxletters}{22}

\DeclareSymbolFont{Symbols}{OMS}{zplm}{m}{n}% Palatino
\DeclareMathSymbol{\infty}{\mathord}{Symbols}{"31}

\usepackage{lscape}
\usepackage{amsthm,tabls}
\usepackage{midpage}
\usepackage{subfigure}
\usepackage{changepage}
\usepackage{longtable}
\usepackage{graphicx,xspace}
\usepackage[a4paper,left=3.54cm,right=2.54cm,top=2.54cm,bottom=2.54cm,centering]{geometry}
\usepackage{makecell}
\newcolumntype{M}[1]{>{\centering\arraybackslash}m{#1}}
\usepackage{multirow}
%\usepackage{showframe}
\usepackage{bigstrut}
\renewcommand\theadfont{\normalsize}
\usepackage{csquotes}
\usepackage{titlesec}
\usepackage{enumitem}
\usepackage{booktabs}
\usepackage{setspace}
\usepackage{multirow}
\usepackage{rotating}

\usepackage{fancyhdr}
\pagestyle{fancyplain}
\fancyhf{}
\fancyhead[R]{\thepage}
\renewcommand{\headrulewidth}{0pt}

\usepackage[apaciteclassic, numberedbib, nosectionbib, tocbib]{apacite}
\usepackage{usebib}
\bibinput{une}
%\printbibliography[heading=none] with apacite
\makeatletter
\let\st@rtbibsection\@bibnewpage
\let\st@rtbibchapter\@bibnewpage
\makeatother
%\printbibliography[heading=none] with apacite

\usepackage[tablewithout,belowskip=-15pt,aboveskip=0pt]{caption}
\renewcommand{\thefigure}{\arabic{figure}}
\renewcommand{\thetable}{\arabic{table}}
\usepackage{url}
\newcommand{\ce}{\centering}
\def\UrlBreaks{\do\.\do\@\do\\\do\/\do\!\do\%\do\7\do\D\do\_\do\\do\;\do\>\do\]%
  \do\)\do\,\do\?\do\'\do+\do\=\do\#}%
\def\UrlBigBreaks{\do\:\do@url@hyp}%
\def\UrlBreaks{\do\/\do-}
%\DefineBibliographyStrings{spanish}{andothers={et~al\adddot}}

%\renewcommand{\thesection}{\Roman{section}}
\renewcommand{\thesubsection}{\arabic{section}.\arabic{subsection}}

\titleformat*{\section}{\normalsize\bfseries}
%\setcounter{secnumdepth}{5}
\newcommand\Tstrut{\rule{0pt}{2.35ex}}
\newcommand\Bstrut{\rule[-1.3ex]{0pt}{1.3ex}}
\newcommand{\TBstrut}{\Tstrut\Bstrut}

\renewcommand{\qedsymbol}{$\blacksquare$}

\newcommand{\ww}{0.015em}
\titlespacing*{\chapter}
{0pt} {5em plus \ww minus \ww} {0em plus \ww minus \ww}
\titlespacing*{\section}
{0pt} {0em plus \ww minus \ww}{0em plus \ww minus \ww}
\titlespacing*{\subsection}
{0pt} {0em plus \ww minus \ww}{0em plus \ww minus \ww} 
\titlespacing*{\subsubsection}
{0pt} {0em plus \ww minus \ww}{0em plus \ww minus \ww}
\titlespacing*{\paragraph} {0pt}{1.25ex plus 1ex minus .2ex}{2em}
\titlespacing*{\subparagraph} {\parindent}{3.25ex plus 1ex minus .2ex}{1em}
\titleformat*{\section}{\normalsize\bfseries}
\titleformat*{\subsection}{\normalsize\bfseries}
\titleformat*{\subsubsection}{\normalsize\bfseries}
\setlength{\parskip}{0.5em plus \ww minus \ww}

\renewcommand{\quote}{\list{}{\rightmargin=\leftmargin\topsep=0pt plus \ww minus \ww}\item\relax}
%%%%%%%%%%%%%%%%%%%%%%%%%%
\usepackage{xpatch}%space equation
\xapptocmd\normalsize{%
	\abovedisplayskip=0.5em plus \ww minus \ww
	\abovedisplayshortskip=0.5em plus \ww minus \ww
	\belowdisplayskip=0.5em plus \ww minus \ww
	\belowdisplayshortskip=0.5em plus \ww minus \ww
}{}{}
%%%%%%%%%%%%%%%%%%%%%
%\setlength\intextsep{0.5em}% plus \ww minus \ww}
\setlist{topsep=0em plus \ww minus \ww, parsep=0.5em plus \ww minus \ww, itemsep=0em} 

\setcounter{topnumber}{2}
\setcounter{bottomnumber}{2}
\setcounter{totalnumber}{4}
\renewcommand{\topfraction}{0.85}
\renewcommand{\bottomfraction}{0.85}
\renewcommand{\textfraction}{0.5}
\renewcommand{\floatpagefraction}{0.7}

\renewcommand{\baselinestretch}{1.2}
%%%%%%%%%%%% flex intertext
\makeatletter
\patchcmd\set@fontsize
{\f@linespread\baselineskip}
{\f@linespread\baselineskip  plus \ww minus \ww}
{}{}
\makeatother

\usepackage{tikz}

\DeclareCaptionLabelSeparator*{spaced}{\\[1.7ex plus .01em minus .01em]}
\captionsetup[table]{textfont=bf,format=plain,labelfont=bf,justification=justified,singlelinecheck=false,labelsep=spaced,skip=0pt}
\captionsetup[figure]{labelfont=bf,justification   = RaggedRight,singlelinecheck = off,skip=5pt,labelfont=it}
%\STsetdecimalsep{{.}}\spanishdecimal{.}

%\setlength\intextsep{0.5em}% plus \ww minus \ww}
\setlength{\textfloatsep}{0.0em plus \ww minus \ww}
%https://tex.stackexchange.com/questions/39017/how-to-influence-the-position-of-float-environments-like-figure-and-table-in-lat


\usepackage[acronym,symbols,nogroupskip]{glossaries-extra}
\usepackage{glossary-superragged}
\makeglossaries

%\usepackage{dsfont}
%\usepackage{textcase}



\usepackage{textcase}
\newcommand{\autor}{Bach. Gutiérrez Brigada, AVELINO LUIS\xspace}
\newcommand{\asesor}{Lic. Canchari Soliz, Jose\xspace}

\newcommand{\expert}{\xspace}
\newcommand{\expertt}{\xspace}

\newcommand{\muestra}{30 estudiantes del colegio público \NoCaseChange{L}os \NoCaseChange{L}icenciados \NoCaseChange{A}yacucho\xspace}

\newcommand{\poblacion}{3000 estudiantes del colegio público \NoCaseChange{L}os \NoCaseChange{L}icenciados \NoCaseChange{A}yacucho\xspace}

\newcommand{\lugar}{Colegio público \NoCaseChange{L}os \NoCaseChange{L}icenciados - \NoCaseChange{A}yacucho\xspace}
%TITULO

\newcommand{\titulo}{LA MÁSCARA ANDINA PARA LA MEJORA DE LA CREATIVIDAD EN ESTUDIANTES DEL NIVEL SECUNDARIO  AYACUCHO 2023\xspace}
% VARIABLES
\newcommand{\variablei}{{{Máscara andina}}\xspace}
\newcommand{\variabled}{{{Creatividad}}\xspace} 
% DIMENSIONES
% VI
\newcommand{\dimi}{{{Dioses y deidades}}\xspace}
\newcommand{\dimii}{{{Personajes ancestrales}}\xspace}
\newcommand{\dimiii}{{{Festivas}}\xspace}
\newcommand{\dimiiii}{{{Hermanos del agua}}\xspace}
% VD 
\newcommand{\dimd}{Valoración estética\xspace}
\newcommand{\dimdd}{Pensamiento crítico\xspace}
\newcommand{\dimddd}{Sensibilidad a los problemas\xspace}
\newcommand{\dimdddd}{Elaboración\xspace}
%%%%%%%%%%%%% Metodologia
\newcommand{\enfoque}{\emph{Cuantitativa}\xspace}
\newcommand{\tipo}{\emph{Pre-experimental}\xspace}
\newcommand{\nivel}{\emph{Explicativa}\xspace}
\newcommand{\diseno}{\emph{Experimental}\xspace}

%%%%%%%%%%% INDICADORES

% INDICADORES DE LA VARIABLE INDEPENDIENTE

% Inteligencia espacial
\newcommand{\gb}{Eventos en el hogar\xspace}
\newcommand{\gbb}{Eventos en el trabajo\xspace}
\newcommand{\gbbb}{Eventos en centros de estudios\xspace}
\newcommand{\gbbbb}{Eventos en centro de esparcimiento\xspace}

\newcommand{\gbbbbb}{Personajes ficticios\xspace}
\newcommand{\gbbbbbb}{Entornos ficticios\xspace}
\newcommand{\gbbbbbbb}{Situaciones absurdas\xspace}
\newcommand{\gbbbbbbbb}{Acciones atemporales\xspace}

\newcommand{\gbbbbbbbbb}{Eventos sociales\xspace}
\newcommand{\gbbbbbbbbbb}{Personajes importantes\xspace}
\newcommand{\gbbbbbbbbbbb}{Representación de protesta\xspace}
\newcommand{\gbbbbbbbbbbbb}{Opinión publica\xspace}

\newcommand{\gbbbbbbbbbbbbb}{Temas de actualidad\xspace}
\newcommand{\gbbbbbbbbbbbbbb}{Temas de controversia\xspace}
\newcommand{\gbbbbbbbbbbbbbbb}{Escándalos sociales\xspace}
\newcommand{\gbbbbbbbbbbbbbbbb}{Corrupción social\xspace}

% INDICADORES DE LA VARIABLE DEPENDIENTE

\newcommand{\fb}{Comprende al prójimo\xspace}
\newcommand{\fbb}{Entiende emociones\xspace}
\newcommand{\fbbb}{Atiende su entorno\xspace}
\newcommand{\fbbbb}{Compasión y apoyo\xspace}

\newcommand{\fbbbbb}{Expresa pensamientos claros\xspace}
\newcommand{\fbbbbbb}{Expresa opinión responsable\xspace}
\newcommand{\fbbbbbbb}{Realiza actos autónomos\xspace}
\newcommand{\fbbbbbbbb}{Resolución de conflictos\xspace}

\newcommand{\fbbbbbbbbb}{Enfoque\xspace}
\newcommand{\fbbbbbbbbbb}{Comprensión\xspace}
\newcommand{\fbbbbbbbbbbb}{Retención\xspace}
\newcommand{\fbbbbbbbbbbbb}{Aplicación\xspace}

\newcommand{\fbbbbbbbbbbbbb}{Procesa información\xspace}
\newcommand{\fbbbbbbbbbbbbbb}{Construye ideas\xspace}
\newcommand{\fbbbbbbbbbbbbbbb}{Formula juicios\xspace}
\newcommand{\fbbbbbbbbbbbbbbbb}{Realiza reflexiones\xspace}



% PROBLEMA
\newcommand{\problema}{¿Como influye la máscara andina para la mejora de la creatividad en estudiantes del nivel secundario  de \NoCaseChange{A}yacucho 2023?\xspace}

\newcommand{\problemae}{¿Como influye la máscara andina en la  composición para la mejora de la creatividad en estudiantes del nivel secundario  de ayacucho 2023.?\xspace}

\newcommand{\problemaee}{¿Como influye la mascara andina en la originalidad para la mejora de l creatividad en estudiantes del nivel secundario de ayacucho 2023??\xspace}

\newcommand{\problemaeee}{¿Como influye la mascara andina en la expresividad para la mejora de la creatividad en estudiantes del nivel secundario de ayacucho 2023?\xspace}

\newcommand{\problemaeeee}{¿Como influye la mascara andina en el contenido para la mejora de la creatividad en estudiantes del nivel secundario de ayacucho 2023?\xspace}



% OBJETIVO
\newcommand{\objetivo}{Determinar la influencia de la máscara andina en la mejora de la creatividad en estudiantes del nivel secundario de ayacucho 2023\NoCaseChange{A}yacucho 2023\xspace}

\newcommand{\objetivoe}{Establecer la influencia de máscara andina en la composición para la mejora de la creatividad en estudiantes del nivel secundario de ayacucho 2023\xspace}

\newcommand{\objetivoee}{Demostrar la influencia de  la máscara andina en la originalidad para la mejora de la creatividad en estudiantes del nivel secundario de ayacucho 2023\xspace}

\newcommand{\objetivoeee}{Explicar la influencia de la máscara andina en la expresividad para la mejora de la creatividad en estudiantes del nivel secundario de ayacucho 2023\xspace}

\newcommand{\objetivoeeee}{Describir como influye la mascara andina en el contenido para la mejora de la creatividad en estudiantes del nivel secundario de ayacucho 2023\xspace}



% HIPOTESIS
\newcommand{\hipotesis}{La máscara andina influye en la mejora de la creatividad en estudiantes del nivel secundario de ayacucho 2023\xspace}

\newcommand{\hipotesise}{La mascara andina influye en la composición para la mejora de la creatividad en estudiantes del nivel secundario de ayacucho 2023\xspace}

\newcommand{\hipotesisee}{La mascara andina  influye  en la originalidad para la mejora de la la creatividad en estudiantes del nivel secundario 2023\xspace}

\newcommand{\hipotesiseee}{La mascara andina influye  en la expresividad para la mejora de la creatividad en estudiantes del nivel secundario de ayacucho 2023\xspace}

\newcommand{\hipotesiseeee}{La mascara andina influye  en el contenido para la mejora de la creatividad en estudiantes del nivel secundario de ayacucho 2023\xspace}



\newcommand{\N}{\mathds{N}}
\newcommand{\R}{\mathds{R}}
\newcommand{\CC}{\mathds{C}}
\newcommand{\I}{\mathds{I}}
\newcommand{\Q}{\mathds{Q}}
\newcommand{\X}{\mathds{X}}


\newcommand{\abs}[1]{\left\vert#1\right\vert}
\newcommand{\norm}[1]{\left\|#1\right\|}
\newcommand{\set}[1]{\left\{#1\right\}}
\newcommand{\seq}[1]{\left<#1\right>}
\newcommand{\co}[1]{\left[#1\right]}
\newcommand{\cc}[1]{\left(#1\right)}
\newglossaryentry{asertividad}{name=Asertividad,description={Asertividad se refiere a la capacidad de expresar tus pensamientos, necesidades y sentimientos de manera clara, honesta y respetuosa, sin dañar a los demás ni permitir que te dañen a ti mismo. La asertividad es una habilidad importante tanto en las relaciones personales como en el ámbito laboral, ya que te permite comunicarte de manera efectiva y establecer límites saludables. La comunicación asertiva implica utilizar un lenguaje claro y directo, expresar tus sentimientos de manera honesta y respetuosa, y escuchar activamente a los demás. Al ser asertivo, puedes mejorar tus relaciones, aumentar tu autoconfianza y mejorar tu bienestar emocional.}}

\newglossaryentry{color}{name=Color,description={Es una propiedad visual que se percibe en los objetos debido a la luz que reflejan o emiten. Los colores son generados por la longitud de onda de la luz visible, que se compone de los colores del espectro (rojo, naranja, amarillo, verde, azul, índigo y violeta) y sus combinaciones. El color puede ser percibido de manera subjetiva y está presente en la naturaleza, el arte, la moda, la decoración, entre otros campos.}}

\newglossaryentry{comunicacion}{name=Comunicación,description={La comunicación es un proceso de intercambio de información entre dos o más personas o entidades. Se lleva a cabo a través de diferentes medios, como lenguaje verbal, escrito, gestual, entre otros, y tiene como objetivo transmitir conocimientos, sentimientos, ideas, mensajes, y establecer un vínculo entre los interlocutores. La comunicación es esencial en todo tipo de relaciones, ya sea a nivel personal, laboral o social. Es una herramienta fundamental para solucionar problemas, tomar decisiones y establecer acuerdos. Por tanto, es importante desarrollar habilidades de comunicación efectiva para poder relacionarnos de manera satisfactoria con las personas que nos rodean. La comunicación puede ser efectiva o inefectiva, dependiendo de la forma en que se desarrolla. Para que la comunicación sea efectiva es necesario que exista una retroalimentación adecuada, una comprensión mutua, respeto y empatía entre las partes involucradas.}}

\newglossaryentry{creatividad}{name=Creatividad,description={La creatividad es la capacidad de generar ideas originales y útiles que permiten resolver problemas o satisfacer necesidades de manera innovadora. También se refiere a la habilidad de crear algo nuevo a partir de elementos ya existentes. La creatividad no se limita a las artes y la cultura, sino que se extiende a todas las áreas de la vida, incluyendo la ciencia, la tecnología, la educación, los negocios y el desarrollo humano. La creatividad puede fomentarse a través de la exploración, la experimentación, la curiosidad y el pensamiento flexible.}}

\newglossaryentry{dibujo}{name=Dibujo,description={Es una representación gráfica de una imagen o idea, realizada con distintas técnicas, como lápiz, carboncillo, tinta, acuarela o técnicas digitales. Puede ser una creación artística o técnica, y se utiliza en diferentes ámbitos, como la arquitectura, el diseño gráfico, la ilustración, la ingeniería y las artes visuales. Los dibujos pueden ser abstractos o figurativos, y pueden transmitir emociones, sensaciones y pensamientos tanto del artista como del espectador.}}

\newglossaryentry{recrear}{name=Recrear,description={Volver a crear algo que ya existió en el pasado. Imitar o representar algo o alguien de manera cercana a la realidad. Divertirse o entretenerse realizando actividades recreativas.}}

\newglossaryentry{arte}{name=Arte,description={El arte es una manifestación cultural que se expresa a través de diversas formas y disciplinas artísticas, como la pintura, la escultura, la música, la literatura, la danza, el cine, el teatro, entre otras. A través del arte, los artistas pueden expresar emociones, sentimientos, ideas y reflexiones sobre la vida y la sociedad. El arte puede ser utilizado como una herramienta de comunicación para transmitir mensajes y provocar cambios sociales, políticos o culturales. Además, el arte tiene un impacto importante en la historia y la evolución de las distintas culturas y civilizaciones. El arte también puede ser una fuente de inspiración y entretenimiento para muchas personas. Puede tener un papel importante en la educación y el desarrollo de la creatividad de los individuos, tanto en términos personales como profesionales. En resumen, el arte es una expresión cultural muy valorada en todo el mundo por su capacidad de conmover, inspirar y transformar.}}

\newglossaryentry{caricatura}{name=Caricatura,description={Las caricaturas son una forma de arte que consiste en realizar una representación exagerada y humorística de una persona o situación. Se utilizan técnicas de dibujo para resaltar ciertas características físicas o de personalidad y se suelen utilizar colores vivos y formas exageradas.}}

\newglossaryentry{comico}{name=Cómico,description={El término "cómico" se refiere a alguien o algo que hace reír o provoca risa. Puede referirse a un humorista, comediante o actor que se dedica a hacer reír a la gente a través de su actuación o sus chistes. También puede referirse a una situación o historia divertida que tiene como objetivo principal provocar risa en el público. En general, un cómico es alguien que se caracteriza por su habilidad para generar humor y alegría en los demás.}}

\newglossaryentry{esbozo}{name=Esbozo,description={En español, la palabra "esbozo" se refiere a la acción de crear una representación inicial o preliminar de algo, ya sea una idea, un dibujo o un plan. Un esbozo puede ser el punto de partida para desarrollar una obra completa, como un boceto antes de hacer un dibujo detallado o un bosquejo antes de escribir un ensayo. También se puede utilizar en sentido figurado para referirse a una descripción general o resumida de algo.}}

\newglossaryentry{periodismo}{name=Periodismo,description={El periodismo es una forma de comunicación que consiste en recopilar, verificar y difundir información sobre eventos, acontecimientos y temas de interés público. Los periodistas son responsables de investigar, escribir y presentar noticias de manera objetiva y precisa.}}

\newglossaryentry{social}{name=Social,description={El término "social" puede tener diferentes significados dependiendo del contexto en el que se utilice. En general, se refiere a todo lo relacionado con la sociedad, las relaciones entre las personas y la interacción social.}}

\newglossaryentry{interrelacion}{name=Interrelación,description={La interrelación se refiere a la conexión o relación que existe entre diferentes elementos, fenómenos o individuos. Implica que dos o más cosas están conectadas, afectándose mutuamente o influyéndose entre sí. La interrelación puede ser de diversos tipos y puede darse en distintos ámbitos, como el social, económico, político, natural, etc.}}

\newglossaryentry{conceptos}{name=Conceptos,description={Una unidad cognitiva de significado. Nace como una idea abstracta (es una construcción mental) que permite comprender las experiencias surgidas a partir de la interacción con el entorno y que, finalmente, se verbaliza (se pone en palabras).}}



\newacronym{UNSCH}{UNSCH}{Uinversidad Nacional De San Cristobal De Huamanga}
%\newacronym{IFS}{IFS}{Iterated function system (sitemas de funciones iteradas)}

\newacronym{VI}{V1}{Variable independiente}
\newacronym{VD}{V2}{Variable dependiente}
\newacronym{V}{V}{Variable}

\newacronym{D}{D}{Dimensiones}

\newacronym{M}{M}{Malo}
\newacronym{B}{B}{Bueno}
\newacronym{R}{R}{Regular}
\newacronym{E}{E}{Excelente}
\newacronym{O}{O}{Observación}

\newacronym{IT}{IT}{$i$--ésimo items de los indicadores de la variable independiente}
\newacronym{ID}{ID}{Identificador de numeración}
\newacronym{Ci}{Ci}{$i$--ésimo items de la variable independiente}
\newacronym{Pi}{Pi}{$i$--ésimo items de la variable dependiente}

%%%%%%%%%%%%%%%%%%%%%%%%%%

%\glsxtrnewsymbol[description={Promedio de las observaciones correspondientes a la dimencion $i$ de la variable dependiente, donde $j=1,2,3,4,5$, $i=1,2,3$ son los promedios en cada observacion $O^i$ de las dimension $Di$ con el método tradicional y experimental respectivamente}]{w1}{$O^i_{\overline{x}_{Dj}^T}$ y $O^i_{\overline{x}_{Dj}^E}$}

\glsxtrnewsymbol[description={Sumatoria de las varianzas de cada elemento muestral}]{var}{$\sum s^2_i$}

\glsxtrnewsymbol[description={Varianzas de cada item}]{var2}{$\sum s^2_i$}

%\glsxtrnewsymbol[description={Promedio total}]{pi1}{$\overline{x}_T$}

%\glsxtrnewsymbol[description={Calificaciones de la variable dependiente en la experimentación tradicional y experimental respectivamente}]{pi2}{$\overline{V}^T$ y $\overline{V}^E$}

%\glsxtrnewsymbol[description={Frecuencia absoluta simple}]{f}{$f$}

%\glsxtrnewsymbol[description={Promedio de los items correspodientes a la dimensión de la variable dependiente con el método tradicional y experimental respectivamente}]{xD1}{$\overline{x}_{Di}^T$ y $\overline{x}_{Di}^E$}

%\glsxtrnewsymbol[description={Promedio de los items correspodientes a la variable dependiente	con el método tradicional y experimental respectivamente}]{xD2}{$\overline{x}_{T}^T$ y $\overline{x}_{T}^E$}


\usepackage{catchfilebetweentags}
\makeatletter
\def\CatchFBT@sanitize{%
   \@sanitize
   \@makeother\{%
   \@makeother\}%
%   \endlinechar=`\^^J% <--- This line modified
}% \CatchFBT@sanitize
\makeatother



\glsaddall

%\usepackage{tocstyle}
%\usetocstyle{standard}

\usepackage{etoolbox}
\makeatletter
\pretocmd{\section}{\addtocontents{toc}{\protect\addvspace{-0.2em plus \ww minus \ww}}}{}{}
\pretocmd{\subsection}{\addtocontents{toc}{\protect\addvspace{0.5em plus \ww minus \ww}}}{}{}
\pretocmd{\subsubsection}{\addtocontents{toc}{\protect\addvspace{0.5em plus \ww minus \ww}}}{}{}
\makeatother

\usepackage{pgfplots}
\usetikzlibrary{intersections}
\pgfplotsset{compat=1.15}
\usepackage{amsmath,pgfplots,gincltex}
\raggedright
\begin{document}
%\SweaveOpts{concordance=TRUE}
\setlength{\abovecaptionskip}{0pt}
\setlength{\belowcaptionskip}{0pt}

%\pagenumbering{roman}
\thispagestyle{empty}

{
\normalsize\centering\thispagestyle{empty}
MINISTERIO DE EDUCACIÓN
\\SUPERINTENDENCIA NACIONAL DE EDUCACIÓN SUPERIOR UNIVERSITARIA (SUNEDU)\\
\vspace{0.65cm}
{\large ESCUELA SUPERIOR DE FORMACIÓN ARTÍSTICA PÚBLICA\\ ``FELIPE GUAMÁN POMA DE AYALA''}\\
{\scriptsize CON NIVEL UNIVERSITARIO: LEY 30220}\\
\vspace{0.65cm}
\includegraphics[height=6.5cm]{logo}\\
\vspace{0.65cm}
{\large TÍTULO}\\
\vspace{0.65cm}
\MakeUppercase{\Large\titulo}\\
\vspace{0.65cm}
{
	\large PRESENTADO POR\\
	\vspace{0.65cm}
	\autor
}\\
\vspace{0.65cm}
\large TPROYECTO DE INVESTIGACIÓN PARA OPTAR EL TÍTULO PROFESIONAL DE LICENCIADO EN EDUCACIÓN ARTÍSTICA\\
\vspace{0.3cm}
ESPECIALIDAD DE ARTES PLÁSTICAS\\
\vspace{0.3cm}
AYACUCHO -- PERÚ\\
\vfill
{2023}

}

\newpage{}
\newpage

\flushbottom

%\renewcommand{\contentsname}{ÍNDICE DE CONTENIDOS}
%%\addcontentsline{toc}{section}{Índice de contenidos}
%\tableofcontents

%\renewcommand\listfigurename{Lista de figuras} 
%\addcontentsline{toc}{section}{Lista de figuras}
%\listoffigures

%\renewcommand\listtablename{Lista de tablas} 
%\addcontentsline{toc}{section}{Lista de tablas}
%\listoftables

%\pagenumbering{arabic}
%\setcounter{page}{1}

%\onehalfspacing
%\begin{center}
%\Large Anteproyecto de Tesis
%\end{center}

%\tableofcontents
\section{DATOS GENERALES}
\subsection{Título} \titulo.
\subsection{Tipo de investigación} Investigación aplicada
\subsection{Autor} \autor
%\subsection{Especialidad} Matemática
\subsection{Asesor} \asesor

\section{PLANTEAMIENTO DEL PROBLEMA}
\subsection{Determinación del problema}




La creatividad a nivel mundial es un fenómeno fascinante en constante evolución. A medida que las fronteras culturales se desdibujan, la tecnología avanza y las personas se conectan en plataformas globales, la creatividad ha adquirido un nuevo alcance y trascendencia.

En todas las disciplinas artísticas y científicas, la creatividad está en constante expansión. Artistas de todo el mundo están explorando nuevas formas de expresión, fusionando diferentes estilos y rompiendo barreras culturales. Esto ha llevado a la creación de obras de arte innovadoras y emocionantes que desafían las convenciones tradicionales.

En el ámbito científico, la creatividad también está en auge. Los investigadores están colaborando a nivel internacional para encontrar soluciones a los desafíos globales, como el cambio climático y las enfermedades, dando lugar a avances significativos en medicina, energía renovable y sostenibilidad.

Además, la creatividad también se ha vuelto fundamental en el mundo empresarial. Las empresas más exitosas son aquellas que son capaces de generar ideas innovadoras y disruptivas, adaptarse rápidamente a los cambios del mercado y crear experiencias únicas para los consumidores. La creatividad se ha convertido en un activo invaluable para la diferenciación y el éxito en un mundo empresarial cada vez más competitivo.

La educación también ha reconocido la importancia de fomentar la creatividad desde temprana edad. Muchos países están implementando reformas educativas que promueven la creatividad y el pensamiento crítico como habilidades fundamentales para el siglo XXI. Esto incluye el fomento de la educación artística y la promoción de entornos de aprendizaje estimulantes que permitan a los estudiantes explorar, experimentar y crear.

En resumen, la creatividad a nivel mundial está en constante cambio y evolución. Es un fenómeno que trasciende fronteras y culturas, impulsando avances en todas las disciplinas y siendo fundamental en la vida moderna. La capacidad de generar ideas nuevas e innovadoras es cada vez más valorada y reconocida como un motor para el progreso y el éxito en el mundo globalizado en el que vivimos.


La creatividad a nivel nacional en Perú es cada vez más reconocida y valorada. El país cuenta con una rica tradición cultural e histórica que ha influido en diversas manifestaciones artísticas y creativas.

En el ámbito de las artes plásticas, tenemos destacados artistas como Fernando de Szyszlo, Teresa Burga, y Juan Javier Salazar, entre otros, cuyas obras han sido exhibidas en museos y galerías de renombre tanto en Perú como en el extranjero.

En la literatura, destacan escritores peruanos como Mario Vargas Llosa, ganador del Premio Nobel de Literatura, y Alfredo Bryce Echenique, cuyas obras han sido traducidas a varios idiomas y han sido reconocidas a nivel internacional.

En la música, la creatividad peruana se ve reflejada en géneros como el criollo, el afroperuano y el huayno, que han evolucionado y se han fusionado con otros ritmos para crear propuestas modernas y originales. Grupos y solistas como Eva Ayllón, Susana Baca y Novalima han logrado trascender las fronteras y llevar la música peruana a distintos escenarios internacionales.

En el ámbito del cine, el Perú ha experimentado un fuerte crecimiento en los últimos años, con películas como "La teta asustada" de Claudia Llosa, que ganó el Oso de Oro en el Festival Internacional de Cine de Berlín en 2009, y "Wiñaypacha" de Óscar Catacora, que fue seleccionada para representar a Perú en los Premios Óscar en 2019.

Además, en el ámbito de la gastronomía, Perú se ha consolidado como uno de los destinos culinarios más destacados a nivel mundial. La creatividad culinaria peruana ha logrado fusionar la diversidad de ingredientes y sabores del país, dando origen a platos reconocidos internacionalmente como el ceviche, el lomo saltado y el ají de gallina.

En conclusión, la creatividad a nivel nacional en Perú se expresa en diversas disciplinas artísticas y culturales, en las que los artistas peruanos han logrado destacar y llevar el talento peruano más allá de sus fronteras.

La creatividad en Ayacucho es un tema muy importante y relevante debido a la rica tradición cultural y artística que existe en esta región de Perú. Ayacucho es conocida por ser cuna de importantes manifestaciones artísticas como la música, la danza, la cerámica y la literatura.

En cuanto a la música, Ayacucho es reconocida por sus festivales folclóricos y por ser la cuna de géneros musicales como el huayno y la música serrana. Los músicos ayacuchanos destacan por su virtuosismo y por la capacidad de fusionar elementos tradicionales con influencias contemporáneas.

En el ámbito de la danza, Ayacucho es famosa por sus coloridos y enérgicos bailes tradicionales como la "Danza de las tijeras" y la "Wititi". Estas danzas son manifestaciones artísticas que combinan ritmo y movimiento de manera única y original.

La cerámica es otra expresión de la creatividad en Ayacucho. Los alfareros ayacuchanos son conocidos por su habilidad en la elaboración de piezas de cerámica adornadas con motivos tradicionales y símbolos ancestrales. Esta actividad artesanal es muy importante en la economía local y se ha convertido en un atractivo turístico.

Por último, la literatura también tiene un lugar destacado en Ayacucho. La ciudad ha sido cuna de importantes escritores peruanos como Eleodoro Vargas Vicuña y Ciro Alegría. La tradición literaria de Ayacucho se ha mantenido viva a lo largo de los años y ha servido como fuente de inspiración para nuevos escritores.

En resumen, la creatividad en Ayacucho se encuentra presente en distintas manifestaciones artísticas y culturales como la música, la danza, la cerámica y la literatura. Estas expresiones creativas son muestra del talento y la riqueza cultural de esta región peruana.

Las máscaras en la creatividad pueden ser una forma de liberación y expresión artística. Al ponernos una máscara, podemos experimentar la libertad para ser alguien diferente, explorar diferentes emociones y roles, y dejar de lado nuestras inhibiciones y miedos.

Además, las máscaras también pueden ser herramientas de juego y experimentación en el proceso creativo. Nos permiten escapar de nuestra identidad y trascender los límites impuestos por la sociedad y nuestras propias mentes. Al ocultar nuestra cara detrás de una máscara, podemos sentirnos más libres y protegidos para explorar nuevas ideas, perspectivas y formas de expresión.

En el arte, las máscaras son utilizadas en diferentes disciplinas creativas, como el teatro, la danza, la pintura y la escultura. Pueden ser utilizadas para representar personajes o emociones específicas, crear narrativas o simplemente como símbolos visuales de la imaginación y la creatividad.

En resumen, las máscaras en la creatividad son herramientas poderosas que nos permiten explorar nuevas facetas de nosotros mismos, liberar nuestras emociones y expresarnos de formas más auténticas. Son un medio para expandir nuestra imaginación y encontrar nuevas formas de expresión artística.

%\begin{displayquote}
%	\emph{Las prácticas semióticas tienen como base los elementos críticos de una
%		actividad semiótica: elección del registro en el cual representar un objeto
%		matemático, capacidad de operar con las transformaciones de tratamiento
%		y de conversión. Dar un sentido coherente a cada representación semiótica
%		de un objeto matemático dado} (p.~57)
%\end{displayquote}




























\subsection{Formulación del problema}
\subsubsection{Problema General}
\problema.
\subsubsection{Problemas Específicos}
\begin{itemize}
	\item \problemae.
	\item \problemaee.
	\item \problemaeee.
	\item \problemaeeee.
\end{itemize}

%\subsection{Importancia y alcances de la investigación}
%\ExecuteMetaData[contenido.tex]{importancia}

\subsection{Justificación y viabilidad}

\subsubsection{Justificación legal}

\subsubsection{Justificación científica}

\subsubsection{Justificación social}

\subsubsection{Justificación metodológica}

\subsubsection{Justificación práctica}


%\subsection{Limitaciones de la investigación}
%\ExecuteMetaData[contenido.tex]{justificacionteo}
%\subsubsection{Limitación metodológica}
%\ExecuteMetaData[contenido.tex]{justificacionteo}
%\subsubsection{Limitación teorica}
%\ExecuteMetaData[contenido.tex]{justificacionteo}
%\subsubsection{Limitación de recurso}
%Los recursos económicos o presupuesto serán totalmente autofinanciados.
%\subsubsection{Limitación temporal}
%\ExecuteMetaData[contenido.tex]{justificacionteo}

\section{OBJECTIVOS}

\subsubsection{Objetivo general}
\objetivo.
\subsubsection{Objetivos específicos}
\begin{itemize}
	\item \objetivoe.
	\item \objetivoee.
	\item \objetivoeee.
	\item \objetivoeeee.
\end{itemize}


\section{MARCO TEÓRICO}

\subsection{Antecedentes de estudio}
\subsubsection{Internacionales}


\subsubsection{Nacionales}

\subsubsection{Regionales}

\subsection{Bases teóricas}
\subsubsection{Contexto de la investigación}

\subsubsection{\variablei}%Representación de objetos matemáticos

Las máscaras andinas son una parte importante de la cultura y la tradición de los pueblos indígenas de los Andes. Estas máscaras se utilizan en diversas festividades, rituales y ceremonias, y cada una tiene un significado simbólico y cultural único. A continuación, se presentan algunos tipos comunes de máscaras andinas:
	Máscaras de animales: representan diferentes animales, como el cóndor, el puma, la serpiente y el zorro. Estas máscaras tienen un gran valor simbólico y espiritual, ya que los andinos consideran a los animales como seres sagrados y guardianes de la naturaleza.

	Máscaras de dioses y deidades: representan a diferentes dioses y deidades de la cosmovisión andina, como Inti (dios del sol), Pachamama (madre Tierra) y Wiracocha (dios creador). Estas máscaras son utilizadas en rituales religiosos y festividades para venerar a estas entidades divinas.

	Máscaras de personajes ancestrales: %representan a los antepasados ​​y líderes históricos de las comunidades andinas. Estas máscaras suelen tener rasgos faciales específicos y están asociadas con cuentos, leyendas y mitos transmitidos de generación en generación.

	Máscaras festivas: se utilizan en festividades y celebraciones tradicionales andinas, como el Carnaval de Oruro en Bolivia y la Fiesta de la Candelaria en Perú. Estas máscaras suelen ser coloridas y elaboradas, y se usan en danzas y procesiones para representar personajes festivos y alegres.

	Máscaras de los hermanos del agua: son máscaras utilizadas en rituales de invocación a las deidades del agua, como los ríos y lagos. Estas máscaras representan seres acuáticos y suelen tener características distintivas, como colores brillantes y diseños de escamas.

	      Estos son solo algunos ejemplos de los tipos de máscaras andinas que existen. Cada comunidad y región puede tener sus propios estilos y diseños únicos de máscaras, pero todas comparten la importancia cultural y espiritual que estas representan.


\subsubsection{\variabled}. Según %\cite{canas_comunicacion_2019}

Existen múltiples dimensiones de la creatividad que han sido identificadas por investigadores y psicólogos. Algunas de estas dimensiones incluyen:

Originalidad: La capacidad de generar ideas nuevas, novedosas y diferentes a lo existente.

Flexibilidad: La habilidad para adaptarse a diferentes situaciones, cambiar de perspectiva y buscar soluciones creativas.

Fluidez: La capacidad de generar una gran cantidad de ideas en un corto período de tiempo.

Elaboración: La habilidad para desarrollar y ampliar ideas, así como para trabajar en los detalles y la implementación.

Sensibilidad a los problemas: La capacidad de identificar y reconocer problemas o situaciones que requieren una solución creativa.

Divergencia: La habilidad para pensar de manera original y fuera de los límites establecidos, explorando diferentes posibilidades y perspectivas.

Valoración estética: La capacidad de apreciar y reconocer la belleza, la originalidad y la creatividad en diferentes formas de expresión.

Pensamiento crítico: La habilidad para cuestionar, analizar y evaluar ideas creativas de manera objetiva.

Estas dimensiones pueden variar en función de la teoría de la creatividad que se utilice, pero en general, ofrecen una visión amplia de los diferentes aspectos que intervienen en el proceso creativo.

%\noindent\textbf{\dimdddd}. De acuerdo a \cite{canas_comunicacion_2019}




\setglossarysection{subsection}\flushbottom \printglossary[numberedsection,nonumberlist,title=Definiciones  de términos básicos]




\section{HIPÓTESIS}
\subsection{Hipótesis}
\subsubsection{Hipótesis principal}
\hipotesis.
\subsubsection{Hipótesis específicas}
\begin{itemize}
	\item \hipotesise.
	\item \hipotesisee.
	\item \hipotesiseee.
	\item \hipotesiseeee.
\end{itemize}

\section{DEFINICIÓN CONCEPTUAL Y OPERACIONALIZACÍÓN DE VARIABLES }
%\subsubsection{Operacionalización de variables}
\begin{itemize}
	\item \textbf{Variable Independiente} (\variablei)
	\item \textbf{Variable Dependiente} (\variabled)
	\item \textbf{Variable Interviniente} (Capacidad artística, interrelación social)
\end{itemize}

\subsubsection{Definición conceptual de variables}


\begin{itemize}
	\item \textbf{Variable Independiente (\variablei):}

	\item \textbf{Variable Dependiente (\variabled):}

	\item  \textbf{Variable Interviniente (Capacidad artística, interrelación social):}


\end{itemize}

\subsubsection{Operacionalización de variables}
\textbf{\variablei}. Para analizar las características del proceso de aplicación de la \MakeTextLowercase{\variablei}, se establecen como dimensiones la \emph{\MakeTextLowercase{\dimi}} y \emph{\MakeTextLowercase{\dimii}}, que cuentan con un total de 10 indicadores, que permitirán determinar el objetivo general esperado por estos procesos, por lo que se aplicara módulos de clases, talleres, exposición y discusión para mejorar el \MakeTextLowercase{\variabled}. Refiérase al Cuadro \ref{www}.

\textbf{\variabled}. El presente trabajo de investigación, se refiere a la habilidad
de incrementar el \MakeTextLowercase{\variabled} lo cual involucra procesos o indicadores como la \MakeTextLowercase{\dimd}, \MakeTextLowercase{\dimdd} y  \MakeTextLowercase{\dimddd}. Se recogerá los datos usando la lista de cotejo, la fichas de observación y pruebas escritas. Para determinar el desarrollo de el \MakeTextLowercase{\variabled}, que cuentan con un total de 16 indicadores.
El cumplimiento determinará el crecimiento de la  eficacia del conocimiento matemático, refiérase al Cuadro \ref{www}.


\begin{table}[ht!]
	\caption{Definición operacional de las variables}
	\label{www}
	\begin{tabular}{|c|l|l|c|c|}
		\hline
		\textbf{{VD}} & \textbf{{D}}                                   & \textbf{Indicadores}             & \textbf{Escala}                                                              & \textbf{Valores}                                                  \\\hline
		\multirow{16}{*}{\rotatebox[origin=c]{90}{\makecell[{{p{5cm}}}]{\centering\variablei}}}
		              & \multirow{3}{*}{\makecell[{{p{3cm}}}]{\dimi}}  & \multirow{1}{*}{\gb}             & \multirow{16}{*}{\rotatebox[origin=c]{90}{\makecell{Cuantitativa de razón}}} & \multirow{16}{*}{\rotatebox[origin=c]{90}{\makecell{Cuaantitaitva \\0 a 20}}} \\\cline{3-3}
		              &                                                & \multirow{1}{*}{\gbb}            &                                                                              &                                                                   \\\cline{3-3}
		              &                                                & \multirow{1}{*}{\gbbb}           &                                                                              &                                                                   \\\cline{3-3}
		              &                                                & \multirow{1}{*}{\gbbb}           &                                                                              &                                                                   \\\cline{2-3}
		              & \multirow{4}*{\makecell[{{p{3cm}}}]{\dimii}}   & \multirow{1}{*}{\gbbbb}          &                                                                              &                                                                   \\\cline{3-3}
		              &                                                & \multirow{1}{*}{\gbbbbb}         &                                                                              &                                                                   \\\cline{3-3}
		              &                                                & \multirow{1}{*}{\gbbbbbb}        &                                                                              &                                                                   \\\cline{3-3}
		              &                                                & \multirow{1}{*}{\gbbbbbbb}       &                                                                              &                                                                   \\\cline{2-3}
		              & \multirow{4}*{\makecell[{{p{3cm}}}]{\dimiii}}  & \multirow{1}{*}{\gbbbbbbbbb}     &                                                                              &                                                                   \\\cline{3-3}
		              &                                                & \multirow{1}{*}{\gbbb}           &                                                                              &                                                                   \\\cline{3-3}
		              &                                                & \multirow{1}{*}{\gbbbbbbbbb}     &                                                                              &                                                                   \\\cline{3-3}
		              &                                                & \multirow{1}{*}{\gbbbbbbbbbbb}   &                                                                              &                                                                   \\\cline{2-3}
		              & \multirow{3}*{\makecell[{{p{3cm}}}]{\dimiiii}} & \multirow{1}{*}{\gbbbbbbbbbbbb}  &                                                                              &                                                                   \\\cline{3-3}
		              &                                                & \multirow{1}{*}{\gbbb}           &                                                                              &                                                                   \\\cline{3-3}
		              &                                                & \multirow{1}{*}{\gbbb}           &                                                                              &                                                                   \\\cline{3-3}
		              &                                                & \multirow{1}{*}{\gbbbbbbbbbbbbb} &                                                                              &                                                                   \\\cline{3-3}
		\hline
		\hline
\multirow{16}{*}{\rotatebox[origin=c]{90}{\makecell[{{p{3cm}}}]{\ce\variabled}}}
& \multirow{4}{*}{\makecell[{{p{2.5cm}}}]{\dimd}}   & \multirow{1}{*}{\fb}               & \multirow{16}{*}{\rotatebox[origin=c]{90}{\makecell{Cuantitativa de proporcioón}}} & \multirow{16}{*}{\rotatebox[origin=c]{90}{\makecell{Cuantitaitva de 0 a 20}}} \\\cline{3-3}
&                                                   & \multirow{1}{*}{\fbb}              &                                                                                    &                                                                               \\\cline{3-3}
&                                                   & \multirow{1}{*}{\fbbb}             &                                                                                    &                                                                               \\\cline{3-3}
&                                                   & \multirow{1}{*}{\fbbb}             &                                                                                    &                                                                               \\\cline{2-3}
& \multirow{4}{*}{\makecell[{{p{2.5cm}}}]{\dimdd}}  & \multirow{1}{*}{\fbbbb}            &                                                                                    &                                                                               \\\cline{3-3}
&                                                   & \multirow{1}{*}{\fbbbbb}           &                                                                                    &                                                                               \\\cline{3-3}
&                                                   & \multirow{1}{*}{\fbbbbbb}          &                                                                                    &                                                                               \\\cline{3-3}
&                                                   & \multirow{1}{*}{\fbbbbbbb}         &                                                                                    &                                                                               \\\cline{2-3}
& \multirow{4}{*}{\makecell[{{p{2.5cm}}}]{\dimddd}} & \multirow{1}{*}{\fbbbbbbbb}        &                                                                                    &                                                                               \\\cline{3-3}
&                                                   & \multirow{1}{*}{\fbbbbbbbbb}       &                                                                                    &                                                                               \\\cline{3-3}
&                                                   & \multirow{1}{*}{\fbbbbbbbbbb}      &                                                                                    &                                                                               \\\cline{3-3}
&                                                   & \multirow{1}{*}{\fbbbbbbbbbbb}     &                                                                                    &                                                                               \\\cline{2-3}
& \multirow{4}{*}{\makecell[{{p{2.5cm}}}]{\dimddd}} & \multirow{1}{*}{\fbbbbbbbbbbbb}    &                                                                                    &                                                                               \\\cline{3-3}
&                                                   & \multirow{1}{*}{\fbbbbbbbbbbbbb}   &                                                                                    &                                                                               \\\cline{3-3}
&                                                   & \multirow{1}{*}{\fbbbbbbbbbbbbbb}  &                                                                                    &                                                                               \\\cline{3-3}
&                                                   & \multirow{1}{*}{\fbbbbbbbbbbbbbbb} &                                                                                    &                                                                               \\\hline

	\end{tabular}
\end{table}


\section{ASPECTO METODOLÓGICO}

\subsection{Enfoque de la investigación}
%Investigación de enfoque cuantitativo debido a que se manipulan datos numéricos con el objetivo de deducir la veracidad de las hipótesis descritas en esta investigación mediante el uso adecuado de las estadísticas descriptiva e inferencial.
%
%De acuerdo a \cite{alzina_metodologiinvestigacion_2004} La orientación cuantitativa se fundamenta en el cálculo de características de los fenómenos sociales, lo cual presupone derivar de un cuadro conceptual adecuado al problema examinado, una cadena de proposiciones que enuncien relaciones entre si las variables experimentadas de forma justificada. Este procedimiento tiende a extender y sistematizar resultados. Según \cite{hernandez_sampieri_metodologiinvestigacion_2014}.

%\begin{displayquote}
%	\emph{La orientación cuantitativa se constituye, de un conjunto de términos que es secuencial y demostrativo. Cada fase antecede a la subsiguiente y no es posible omitir o evitar sucesos. La disposición es inflexible. Parte de una idea que va delimitándose y, una vez definida, se proceden a generar objetivos y cuestiones de exploración, se revisa la bibliografía y se elabora un marco o una configuración teórica. De las cuestiones se forman hipótesis y establecen variables; se bosqueja un procedimiento para experimentar la cual se conoce como diseño, se calculan las variables en un determinado contexto; se exploran los datos obtenidos, recurriendo a procesos estadísticos, y se deduce una serie de conclusiones con relación a la o las hipótesis}. (p. 97)
%\end{displayquote}

\subsection{Método de investigación}

%La investigación es de tipo aplicada pues esta investigación busca una teoría didáctica que pueda ser aplicada en el campo educativo una vez sea probada la veracidad y objetividad de las hipótesis planteadas.

\subsection{Tipo de investigación}
%La investigación es de tipo aplicada, que tiene por finalidad contribuir al conocimiento científicos de aprendizaje y herramientas pedagógicas. De acuerdo a \cite{alzina_metodologiinvestigacion_2004}.%\citeauthor{Duval} (1999, p. 15 citado en \citeauthor{ines}, \citeyear{ines}, p.~15).
%
%\begin{displayquote}
%	\emph{Es también llamada práctica, empírica, activa y se encuentra íntimamente ligada a la investigación básica, ya que depende de sus descubrimientos y aportes teóricos para poder generar beneficios y bienestar a la sociedad. Se fundamenta en la investigación teórica; su finalidad especifica es aplicar las teorías existentes a la producción de normas y procedimientos tecnológicos para controlar situaciones  o procesos de la realidad.} (p.~15)
%\end{displayquote}

%Al respecto \cite{kothari_research_2004}, describe la investigación aplicada, como el manejo de las sapiencias, en la experiencia cotidiana, para emplearlos, en la totalidad de los casos, en beneficio de la humanidad.
%
%\cite{hernandez_sampieri_metodologiinvestigacion_2014}. ``Se refiere a un estudio de investigación en
%el que se manipulan deliberadamente una o más variables independientes
%(supuestas causas)
%para analizar las consecuencias de esa manipulación sobre una o más
%variables dependientes
%(supuestos efectos), dentro de una situación de control para el investigador''.


\subsection{Diseño de investigación}
El diseño de investigación \MakeTextLowercase{\diseno}, que se estructura de acuerdo a la Tabla \ref{figg}.
%individuos de la población iguales oportunidades de ser seleccionados, porque los grupos ya están formados con los estudiantes de la serie 200 en la \lugar.


\begin{table}[ht!]
	\caption{Diesño de investigación}
	\label{figg}
	\centering
	\begin{tabular}{cccc}
		\toprule
		Pre & Experimento & Post \\
		\midrule
		O1  & X1          & O2   \\
		\bottomrule
	\end{tabular}
\end{table}

%\begin{figure}[ht!]\centering
%	\begin{tikzpicture}[> = stealth,shorten > = 2pt,semithick]
%\node[] at (0,0)  (T){$M$};
%\node[] at (2,1.3)  (o1){$O_1$};
%\node[] at (5,1.3)  (w){$V_1$};
%\node[] at (2,0)  (r){$r$};
%\node[] at (4,1.3)  (o2){$O_2$};
%\node[] at (5,-1.3)  (ww){$V_2$};
%\path[->] (o1) edge node {} (o2);
%\path[->] (T) edge node {} (ww);
%\path[<-] (o1) edge node {} (w);
%\path[<-] (o2) edge node {} (ww);
%\path[->] (o1) edge node {} (r);
%\path[->] (o2) edge node {} (r);
%	\end{tikzpicture}
%	\caption{Diseño \MakeTextLowercase{\diseno}}
%	\label{figg}
%\end{figure}

%Donde $M$ es la muestra, $O_1$ es la variable 1, \MakeTextLowercase{\variablei},
%$O_2$ es la variable 2, \MakeTextLowercase{\variabled} y
%$r$ es la relación entre variable 1 y variable 2.
%
%El diseño de la investigación es cuasi -- experimental de un mismo grupo de trabajo con pre y pos -- prueba, porque se tomará un solo grupo experimental. De acuerdo a textcite{design} este diseño permitirá aplicar los módulos de experimentación en un periodo determinado y luego se trabajará de manera tradicional, alternando sucesivamente.
%
%\begin{displayquote}
%	\emph{Al elegir un diseño se tiene en cuenta una secuencia de procesos para aplicarlos en un determinado problema y deducir resultados que verifiquen las hipótesis del problema. En este caso se tendrá el diseño propuesto con un solo grupo, que permitirá aplicar la primera variable sobre la segunda, para fortalecer la elaboración plástica tridimensional.}. \cite{kothari_research_2004}[p 243]
%\end{displayquote}



\subsection{Población y muestra}

\subsubsection{Población}

%Constituida por \poblacion, matriculados en el 2023. Los criterios que se considerarán en la selección de la muestra se describen en el Tabla \ref{apt7}.
%
%
%\begin{table}[ht!]
%	\caption{Criterio de inclusión y exclusión}\label{apt7}
%	\centering
%	\begin{tabular}{ccc}\Xhline{2pt}
%		Población                                                                        & Aptos & No aptos \\\midrule
%		\makecell*[{{p{3.5cm}}}]{\centering Estudiantes matriculados del semestre impar} &
%		\makecell*[{{p{3.9cm}}}]{\centering Estudiantes regulares                                           \\ Asistentes puntuales} &\makecell[c]{ Estudiantes repitentes \\
%			Estudiantes retirados                                                                               \\
%			Estudiantes del quinto superior                                                                     \\
%			Estudiantes no asistentes}                                                                          \\\bottomrule
%	\end{tabular}
%\end{table}


%De acuerdo a \cite{zacarias_metodologiinvestigacion_2020}, la población o universo es el agregado de todos los casos que concuerdan con determinadas descripciones o características, es decir es un conjunto de elementos que tiene al menos una misma característica que pueden ser medidos cuantitativamente o cualitativamente.
%
%\subsubsection{Muestra}
%
%En la presente investigación la muestra será no probabilística e intencional compuesta por un solo grupo experimental de \muestra, al cual se le aplicará el diseño \MakeTextLowercase{\diseno} de un mismo grupo con dos tipos de pruebas en series temporales equivalentes (Alternado).
%
%Según \cite{hernandez_sampieri_metodologiinvestigacion_2014}, la muestra es una subcolección de elementos de la población seleccionada, de donde se obtendrá información para el desarrollo del estudio y del cual se hará medicines y observaciones de la variable dependiente e independiente, interrelacionada.

\subsubsection{Muestreo}

%El muestreo será no probabilístico e intencional. Esta técnica consiste  en obtener una muestra donde los elementos  se recogen de manera aleatoria.
%
%\cite{alzina_metodologiinvestigacion_2004} se refiere a la muestra no probabilística como aquella donde las unidades de la población tiene la misma posibilidad de ser elegidas y pertenecer a la muestra. Es deliberado, ya que esta práctica se opera en poblaciones indistintas. Aquí el especialista, conociendo la población y con buen juicio resuelve que características de análisis compondrá la muestra.


\subsection{Técnicas e instrumentos de recolección de datos}

\subsubsection{Técnica de recolección de datos}

%\paragraph{Observación} Para el análisis de la influencia de la variable independiente sobre la variable dependiente se utilizó las técnicas de experimentación y observación, empleándose como instrumentos el plan de experimentación y la ficha de observación. Esta técnica que permitió recoger datos de la influencia de la variable independiente \MakeTextLowercase{\variablei} sobre la variable dependiente \MakeTextLowercase{\variabled}, en el proceso de experimentación.
%
%
%%
%Según \cite{paz_metodologiinvestigacion_2014} el método de observación es el método más utilizado, especialmente en estudios relacionados con ciencias. La observación se convierte en una herramienta científica y en el método de recolección de datos para el investigador, cuando sirve a un propósito de investigación formulado, se planifica y registra sistemáticamente y se somete a controles de validez y fiabilidad.
%
%\paragraph{Prueba pedagógica} Técnica que permitió recoger de manera escrita del  \MakeTextLowercase{\variablei} logro del aprendizaje de la variable dependiente \MakeTextLowercase{\variabled}, en el proceso de la experimentación.
%
%\cite{hernandez_sampieri_metodologiinvestigacion_2014} refiere: ``Es un proceso a través del cual se compara una unidad preestablecida y que la evaluación es un proceso que consiste en obtener infom1ación sistemática y objetiva acerca de un fenómeno e interpretar dicha información a fin de seleccionar entre distintas alternativas de decisión'' (p.~68).
%
%\paragraph{Experimental} Técnica que permitió aplicar los módulos de experimentación de la variable \variablei.
%
%
%{Según \cite{kothari_research_2004} el método de observación es el método más utilizado, especialmente en estudios relacionados con ciencias. La observación se convierte en una herramienta científica y en el método de recolección de datos para el investigador, cuando sirve a un propósito de investigación formulado, se planifica y registra sistemáticamente y se somete a controles de validez y fiabilidad.}

\subsubsection{Instrumento de recolección de datos}

%\paragraph{Plan de experimentación.} Documento pedagógico que contiene el proceso de aplicación de las estrategias pictóricas (ver anexo).
%
%\paragraph{La ficha de observación} En el cual se incluyen indicadores que permitieron conocer  los indicadores de la variable independiente planteadas que influirán en el desarrollo de los indicadores de la variable dependiente en la muestra. En este caso el proceso será observado a través del evento práctico teórico que hicieron los observados.
%


%La ficha de observación por \cite{bernal}, es la exploración visual de lo que sucede en un contexto existente, en un fenómeno explícito, especificando y estableciendo los hechos acertados de acuerdo con algún diseño previsto, manipula un herramienta o prontuario impreso, reservado a conseguir respuestas sobre el problema en estudio.


%
%\paragraph{La ficha de opinión}
%
%Este instrumento permitirá recoger la opinion de los estudiantes con respecto al uso de la \MakeTextLowercase{\variablei}, y averiguar el impacto de esta. De acuerdo a cite{opinion}. Una ficha de opinión, es una ficha en la que personas externas a la investigación escribe plasma lo que piensa en relación al texto o estudio que se está realizando.
%
%
%
%Esta ficha puede tener: Número (es para tener orden de las fichas), Autor (la persona quien lo escribe), Tema de que se trata, Items o cuestionarios y Nota (en caso requerido).
%
%
%
%%URL del artículo: %http://www.ejemplode.com/13-ciencia/2411-ejemplo_de_ficha_de_opinion.html
%%Leer completo: ejemplos de Ficha de opinión
%%Los indicadores que conformaran la estructura de la ficha de
%%observación fueron elaborados a partir de la operacionalización
%%de la variable independiente, siendo las posibilidades de sus
%%respuestas: No, A veces y Sí.
%\paragraph{Prueba escrita} Instrumento que permitirá recoger datos del desarrollo de las capacidades matemáticas  \MakeTextLowercase{\dimd},  \MakeTextLowercase{\dimdd}, \MakeTextLowercase{\dimddd} y \MakeTextLowercase{\dimddd}; instrumento elaborado de acuerdo a los indicadores descritos en el presente trabajo de investigación.
%
%Las competencias y sus valoraciones cualitativas y cuantitativas se describen en la siguiente Tabla \ref{evaluacion}.
%
%\begin{table}[ht!]
%	%\centering
%	\caption{Criterio de calificación}\label{evaluacion}
%	\begin{tabular}{cccc}\Xhline{2pt}
%		\bf \multirow{2}*[-.01cm]{Capacidades}                                                   & \bf Valoración   & \bf Valoración    \\
%		                                                                                         & \bf  cualificada & \bf  cuantificada \\\midrule
%		\multirow{5}*[-.01cm]{\makecell*[{{p{7cm}}}]{\centering  \dimd  \dimdd \dimddd \dimddd}} & Excelente        & 17 -- 20          \\
%		                                                                                         & Bueno            & 13 -- 16          \\
%		                                                                                         & Regular          & 09 -- 12          \\
%		                                                                                         & Malo             & 05 -- 08          \\
%		                                                                                         & Deficiente       & 00 -- 04          \\
%		\bottomrule
%	\end{tabular}
%\end{table}

\subsection{Tratamiento estadístico de los datos}
Con el objetivo de procesar adecuadamente los datos y sacar conclusiones confiables, se utiliza  las siguientes estadísticas.

%\subsubsection{La estadística descriptiva} Esta estadística se utiliza en el estudio  del comportamiento de los datos, para inferir conclusiones relativas de nuestras variables y la relación entre ellas.
%
%\begin{enumerate}
%	\item Tendencia central (media, mediana y moda),
%	\item dispersión o variabilidad (desviación estándar, la varianza y la regresión estándar),
%	\item curtosis (leptocurtica, mesocúrtica y platicúrtica) y
%	\item asimetría (asimétrica positiva, simétrica y asimétrica negativa).
%\end{enumerate}
%
%\subsubsection{La estadística inferencial} Esta estadística se utilizará para probar, la normalidad y la homocedasticidad  de los datos, además de las prueba de hipótesis.
%
%Cabe mencionar que se hará uso del software R Sweave y Excel para el procesamiento de los datos.

\subsubsection{Validación de instrumentos}
Consistirá en validar el contenido del instrumento por el criterio expertos (con maestría o doctorado), quienes verificarán y evaluarán la coherencia y secuencialidad de los instrumentos.



La opinión de los expertos consultados permitirá establecer la validez de los instrumentos que se empleará en la investigación. La validez se establecerá mediante el método del juicio de expertos.



%\begin{adjustwidth}{1cm}{}
%	\emph{La validez, en términos generales, se refiere al grado en que un instrumento mide
%		realmente la variable que pretende medir. Por ejemplo, un instrumento válido para medir la inteligencia debe medir la inteligencia y no la memoria. Un método para medir
%		el rendimiento bursátil tiene que medir precisamente esto y no la imagen de una empresa. Un ejemplo aunque muy obvio de completa invalidez sería intentar medir el peso de los objetos con una cinta métrica en lugar de con una báscula.} \cite[p.~200]{alzina_metodologiinvestigacion_2004}
%\end{adjustwidth}

%\subsection{Procedimientos}
%\subsection{Estadísticos descriptivos}
%\subsection{Estadísticos inferenciales}

\subsubsection{Cofiabilidad de instrumentos}
Se usara el coeficiente de Alpha Chrombach, Refierase a la tabla \ref{pret:1} ...
%Existen  diversos procedimientos para calcular la confiabilidad de un instrumento de medición. Todos utilizan formulas que producen coeficientes de confiabilidad. La mayoría de estos coeficientes pueden oscilar entre cero y uno donde un coeficiente cero significa nula confiabilidad y uno representa  un máximo de confiabilidad.

%La confiabilidad de consistencia interna del instrumento, se determinará con la prueba piloto, en una muestra de 10 estudiantes que no serán miembros de la muestra, aplicando mitades de Coeficiente de Pearson y la corrección de Spearman Brow, la fórmula referencial es la siguiente: $$r_{xy}=\frac{n\Xsum xy-\cc{\Xsum x}\cc{\Xsum y}}{\sqrt{n\Xsum x^2-\cc{\Xsum x}^2}\sqrt{n\Xsum y^2-\cc{\Xsum y}^2}}$$
%
%La corrección con Spearman Brow $$r=\frac{2r_{xy}}{1+r_{xy}}$$
%
%Donde:
%$n$ tamaño de muestra
%$x$ es la puntuación de la primera mitad (Impares)
%$y$ es la puntuación de la segunda mitad (Pares)
%$r_{xy}$ es el coeficiente Pearson
%$r$ es el coeficiente de Spearman Brow. Si $0\leq r\leq 0.8$ se considera no confiable y si $0.8< r \leq1$ se considera confiable.

%\begin{displayquote}
%	\emph{La confiabilidad de un instrumento de medición es el grado en que su aplicación
%		repetida al mismo individuo u objeto genera resultados similares. Por ejemplo, si se midiera en este momento la temperatura ambiental usando un termómetro y éste indicara que hay $22^\circ C$, y un minuto más tarde se consultara otra vez y señalara $5^\circ C$, tres minutos después se observara nuevamente y éste indicara $40^\circ C$, dicho termómetro no sería confiable, ya que su aplicación repetida produce resultados distintos. Asimismo, si una prueba de inteligencia (Intelligence Quotient, IQ) se aplica hoy a un grupo de personas y da ciertos valores de inteligencia, se aplica un mes después y proporciona valores diferentes, al igual que en subsecuentes mediciones, tal prueba no sería confiable. Los resultados no son coherentes, pues no se puede confiar en ellos.
%	} \cite[p.~61]{kothari_research_2004}\end{displayquote}

\begin{table}[ht!]
	% \centering\STautoround{3}
	{%\tiny
		\caption{Prueba de confiabilidad de la ficha de observación}\label{pret:1}
		%\vspace{0.5cm}
		\begin{tabular}{|c|c|c|c|c|c|c|c|c|c|c|c|}\hline
			AI             & \multicolumn{10}{c}{ITEMS} & $\sum$                                                         \\\hline
			               & IN1                        & IN2    & IN3 & IN4 & IN5 & IN6 & IN7 & IN8 & $\ldots$ & IN23 & \\\hline
			1              &                            &        &     &     &     &     &     &     &          &      & \\\hline
			2              &                            &        &     &     &     &     &     &     &          &      & \\\hline
			3              &                            &        &     &     &     &     &     &     &          &      & \\\hline
			$\vdots$       &                            &        &     &     &     &     &     &     &          &      & \\\hline
			20             &                            &        &     &     &     &     &     &     &          &      & \\\hline
			$\sum$         &                            &        &     &     &     &     &     &     &          &      & \\\hline
			$\overline{x}$ &                            &        &     &     &     &     &     &     &          &      & \\\hline
			$s_i^2$        &                            &        &     &     &     &     &     &     &          &      & \\\hline
			$s_i$          &                            &        &     &     &     &     &     &     &          &      & \\\hline
		\end{tabular}\\\vspace{0.5cm}
		{\normalsize Fuente: Datos obtenidos de una prueba piloto}
	}
\end{table}


%El pre test con los datos obtenidos mediante la calificación de los expertos en artes plásticas la ficha de observación generó ... de acuerdo a la Tabla ref{pret}, el instrumento (la ficha de observación)es confiable pues $0.8<\alpha<1$.

%$$\alpha=\frac{k}{k-1}\left(1-\frac{\sum s_i^2}{s_t^2}\right)$$


\section{ASPECTO  ADMINISTRATIVO}
\subsection{Recursos}

\subsubsection{Recursos humanos}
\begin{enumerate}
	\item Responsable de la investigación
	\item Profesor Asesor
	\item Estudiantes de la \lugar.
	\item Autoridades de el \lugar.
	\item Autoridades de la Universidad ESFAPA "FGPA".
\end{enumerate}

\subsubsection{Recursos materiales}
\begin{enumerate}
	\item Bibliografía especializada e Internet %\cite{geer}
	      %\item Materiales y equipos de laboratorio %\cite{ctan}
	\item Materiales de escritorio y computadora %\cite{rahtz:96-1}
	\item Otros materiales %\cite{girou:01:}
\end{enumerate}

\subsubsection{Recursos financieros}
\begin{enumerate}
	\item Auditorio de la \lugar
	\item Laptop y Data Show
\end{enumerate}

\subsection{Presupuesto}
Materiales y equipos (financiados) refiérase al Cuadro \ref{presup}.

\begin{table}[ht!]
	\caption{Presupuesto}\label{presup}
	\begin{tabular}{|c|c|c|}\hline
		$\text N^\circ$ & Materiales y servicios              & {Gastos en S/} \\\hline
		1               & Materiales didácticos               & 135.30         \\\hline
		2               & Impresión                           & 102.50         \\\hline
		3               & Libros                              & 300.90         \\\hline
		4               & Empastados de la tesis (4 unidades) & 80.10          \\\hline
		5               & Gastos de movilidad                 & 90.00          \\\hline
		6               & Otros gastos                        & 747.00         \\\hline
		                & TOTAL                               & 1455.8         \\
		\hline
	\end{tabular}
\end{table}

\subsection{Cronograma de actividades}
Temporalización-cronograma de acciones refiérase al Cuadro \ref{cronog}.
\begin{table}[ht!]

	\caption{Cronograma de acciones}\label{cronog}
	\begin{tabular}{|c|p{6cm}|c|c|c|c|c|c|c|c|c|c|}
		\hline
		\ce $\text N^\circ$ & Actividades                       & \multicolumn{10}{c|}{2023}                                                                                                             \\
		\hline
		                    &                                   & E                          & M         & J         & J         & A         & S         & O         & N         & D         & E         \\\cline{3-12}
		1                   & Análisis bibliográfico            & $\bullet$                  & $\bullet$ &           &           &           &           &           &           &           &           \\
		2                   & Elaboración del marco teórico     & $\bullet$                  & $\bullet$ &           &           &           &           &           &           &           &           \\
		3                   & Elaboración de los instrumentos   &                            & $\bullet$ & $\bullet$ &           &           &           &           &           &           &           \\
		4                   & Prueba de los instrumentos        &                            &           & $\bullet$ &           &           &           &           &           &           &           \\
		5                   & Aprobación del proyecto           &                            &           &           & $\bullet$ & $\bullet$ & $\bullet$ &           &           &           &           \\
		6                   & Recolección de datos              &                            &           &           &           &           &           & $\bullet$ & $\bullet$ &           &           \\
		7                   & Procesamiento y análisis de datos &                            &           &           &           &           &           &           & $\bullet$ &           &           \\
		8                   & Redacción de la tesis             &                            &           &           &           &           &           &           &           & $\bullet$ &           \\
		9                   & Presentación de la tesis          &                            &           &           &           &           &           &           &           & $\bullet$ &           \\
		10                  & Sustentación                      &                            &           &           &           &           &           &           &           &           & $\bullet$ \\
		\hline
	\end{tabular}
\end{table}



\subsection{Potencial humano}
Personas que participaron en la realización de la investigación

\begin{table}[ht]
	\centering
	\caption{Potencial humano}\label{potencialhumano}
	\begin{tabular}{ccc}
		\Xhline{2pt}
		@$\text N^\circ$ & @POTENCIAL HUMANO                          & @CANTIDAD \\\midrule
		1                & Alumnos ($1^\circ$ artistas)               & 12        \\
		2                & Profesores de la especialidad de escultura & 3         \\
		3                & Autodidactas                               & 2         \\
		4                & Estadistas                                 & 3         \\
		5                & Físicos                                    & 2         \\
		6                & Matemáticos                                & 2         \\\midrule
		                 & TOTAL                                      & 7         \\
		\bottomrule
	\end{tabular}
\end{table}



\section{REFERENCIAS}
%\renewcommand{\bibname}{REFERENCIAS BIBLIOGRÁFICAS}
\bibliographystyle{apacite}
\bibliography{une}

\vspace{-0.3cm}

%\setglossarysection{section}
%\setlength{\glsdescwidth}{0.7\textwidth}
%\printglossary[nopostdot,style=long,nogroupskip,nonumberlist,type=\acronymtype,title=ACRÓNIMOS]
%\vspace{-0.3cm}

%\setglossarysection{section}
%\setlength{\glsdescwidth}{0.7\textwidth}
%\printglossary[nopostdot,style=long,nogroupskip,nonumberlist,type=symbols,title=SÍMBOLOS]
%\vspace{-0.3cm}


\section{ANEXOS}
%\addcontentsline{toc}{section}{APÉNDICES}
%\renewcommand{\baselinestretch}{1.0}
%\pagenumbering{Roman}

%\setcounter{subsection}{0}
%\renewcommand{\thesubsection}{\Alph{subsection}}
%%\renewcommand{\thesubsubsection}{A.\arabic{subsection}}

\numberwithin{equation}{subsection}
\numberwithin{figure}{subsection}
\numberwithin{table}{subsection}

\subsection{Ficha técnica\label{ficha tecnica}}
\begin{enumerate}
	\item \textbf{Titulo }\titulo\textbf{ }
	\item \textbf{Año de realización }2019
	\item \textbf{Autor de la investigación }\autor
	\item \textbf{Asesor de la investigación }\asesor
	\item \textbf{Colaboradores que validaron los instrumentos }\asesor, \expert y \expertt.
	\item \textbf{Institución } \lugar
	\item \textbf{Escuela profesional }\lugar
	\item \textbf{Formato }PDF.\textbf{ }

	      %	\item \textbf{Resumen }El presente trabajo de investigación tiene como objetivo 	\MakeTextLowercase{\objetivo}. Tipo de investigación básica, nivel 	de investigación experimental explicativa, de diseño cuasiexperimental 	de un grupo con pre y posprueba en series temporales equivalentes 	(alternado); se empleo el método hipotético deductivo, experimental	y estadístico descriptivo e inferencial; el lugar de estudio fue en	la \lugar; la muestra fue no probabilística e intencional compuesta	por un solo grupo experimental de 35 estudiantes de la serie 200,	matriculados en el curso de arquitectura, del semestre impar de la	escuela profesional de Ingeniería Civil UNSCH; los datos fueron recolectados	a través de la prueba escrita y la ficha de observación; la prueba 	de validez de instrumentos se realizó a través de juicio de expertos 	y la confiabilidad, a través de prueba del Coeficiente de Pearson 	y la corrección de Spearman Brow. Se verificó la no normalidad de 	los datos, mediante la prueba de  \emph{Shapiro Wilks}; se aplicó 	la prueba de Student dos muestras relacionadas para la prueba de hipótesis, 	con un nivel de confianza del 95\% y significancia del 5\% y se concluyó, 	que \MakeTextLowercase{\hipotesis}. \\Palabras Claves: \variablei, \MakeTextLowercase{\upshape \variabled, \dimddd, \dimi y \dimii}. 
	\item \textbf{Estructura de la tesis \flushbottom }
	      \begin{enumerate}
		      \item PLANTEAMIENTO DEL PROBLEMA
		            \begin{enumerate}
			            \item Identificación y descripción del problema
			            \item Formulación del problema
			            \item Objetivos
			            \item Justificación de la investigación
		            \end{enumerate}
		      \item MARCO TEÓRICO
		            \begin{enumerate}
			            \item Antecedentes de la investigación
			            \item Bases teóricas
			            \item Definiciones de términos básicos
		            \end{enumerate}
		      \item METODOLOGÍA DE LA INVESTIGACIÓN
		            \begin{enumerate}
			            \item Sistema de hipótesis
			            \item Sistema de variables
			            \item Operacionalización de variables
			            \item Aspecto metodológico
		            \end{enumerate}
		      \item RESULTADOS DE LA INVESTIGACIÓN
		            \begin{enumerate}
			            \item Análisis e interpretación
			            \item Resultados inferenciales
			            \item Discusión de resultados
		            \end{enumerate}
		            CONCLUSIONES

		            RECOMENDACIONES

		            REFERENCIAS BIBLIOGRÁFICAS

		            ANEXOS
	      \end{enumerate}
	\item \textbf{Antecedentes }

	\item \textbf{Problema }
	      \begin{enumerate}
		      \item \textbf{Problema general.}

		            \problema.
		      \item \textbf{Problemas específicos}
		            \begin{itemize}
			            \item \problemae.
			            \item \problemaee.
			            \item \problemaeee.
			            \item \problemaeeee.
		            \end{itemize}
	      \end{enumerate}
	\item \textbf{Objetivos }
	      \begin{enumerate}
		      \item \textbf{Objetivo general}

		            \objetivo.

		      \item \textbf{Objetivos específicos }
		            \begin{itemize}
			            \item \objetivoe.
			            \item \objetivoee.
			            \item \objetivoeee.
			            \item \objetivoeeee.
		            \end{itemize}
	      \end{enumerate}
	\item \textbf{Hipótesis }
	      \begin{enumerate}
		      \item \textbf{Hipótesis general}

		            \hipotesis.
		      \item \textbf{Hipótesis específicas}
		            \begin{itemize}
			            \item \hipotesise.
			            \item \hipotesisee.
			            \item \hipotesiseee.
			            \item \hipotesiseeee.
		            \end{itemize}
	      \end{enumerate}
	\item \textbf{Orientación metodológica y diseño muestral}
	      \begin{itemize}
		      \item \textbf{Enfoque} Cuantitativo.
		      \item \textbf{Tipo} Investigación aplicada.
		      \item \textbf{Nivel de estudio} Investigación explicativa experimental.
		      \item \textbf{Diseño} Diseño \MakeTextLowercase{\diseno} de un mismo grupo
		            con pre y postprueba en series temporales equivalentes (Alternado).
		      \item \textbf{Metodología} Experimental, hipotético deductivo, estadístico
		            y analítico.
		      \item \textbf{Población} Constituida por \poblacion.
		      \item \textbf{Muestra} Constituida por \muestra.
		      \item \textbf{Muestreo} No probabilístico.
		      \item \textbf{Técnica} Observación, evaluación pedagógica y experimentación.
		      \item \textbf{Instrumentos} Ficha de observación, ficha de opinión prueba
		            escrita y módulos experimentales.
		      \item \textbf{Procesamiento} \textbf{de datos} Mediante los programas estadísticos SPSS y Excel.
	      \end{itemize}
	\item \textbf{Análisis y presentación de resultados} Se procesará los datos
	      recogidos en el salón de clases, es decir, se sistematizará, interpretará
	      y se sintetizará para responder así a los objetivos y a la finalidad
	      de la investigación, con la perspectiva, de que la presentación de
	      los resultados de esta investigación sirva de reflexión, sobre el
	      abordaje de la \MakeTextLowercase{\variablei}, como un factor
	      fomentador del \MakeTextLowercase{\variabled}.
	\item \textbf{Conclusión }\hipotesis\textbf{.}
\end{enumerate}
\textbf{\vspace{0.5cm}
}
\begin{center}
	....................................................\\
	\autor
\end{center}

\clearpage

\begin{landscape}
	\subsection{Matriz de consistencia\label{consistencia}}

	%\lugar -- TÍTULO: \titulo -- AUTOR: \autor

	\begin{table}[ht!]\caption{Matriz de consistencia}
		\centering \scriptsize  \renewcommand\tabcolsep{0.1cm}\renewcommand\arraystretch{1}

		\begin{tabular}{|p{3.9cm}|p{3.9cm}|p{3.9cm}|c|c|c|} \hline \ce\bf PROBLEMAS & \ce\bf OBJETIVOS     & \ce\bf HIPÓTESIS & \ce\bf VARIABLES & \ce\bf MARCO TEÓRICO & \bf METODOLOGÍA\Tstrut \\ \hline \ce \bf PROBLEMA GENERAL &\ce \bf OBJETIVO GENERAL &\ce \bf HIPÓTESIS GENERAL &
			\multirow{8}*[0cm]{
				\begin{minipage}[t]{2.9cm}
					\begin{itemize}[itemsep=-0pt,leftmargin=*,labelsep=.02cm,] \item \textbf{Variable 1}: \variablei.
						      \begin{itemize}[itemsep=-0pt,leftmargin=*,labelsep=.02cm,topsep=-2pt]
							      \item \dimi.
							      \item \dimii
							      \item \dimiii.
							      \item \dimiiii.
						      \end{itemize} \item \textbf{Variable 2}: \variabled.
						      \begin{itemize}[itemsep=-0pt,leftmargin=*,labelsep=.02cm,topsep=-2pt]
							      \item \dimd.
							      \item \dimdd.
							      \item \dimddd.
							      \item \dimdddd.
						      \end{itemize}
					\end{itemize}
				\end{minipage}
			}                                                            & \multirow{8}*[0cm]{
				\begin{minipage}[t]{2.9cm}
					\begin{itemize}[itemsep=-0pt,leftmargin=*,labelsep=.02cm,]
						\item Contexto de la investigación.
						\item \variablei.
						\item \variabled.
						\item \dimi.
						\item \dimii.
						\item \dimiii.
						\item \dimiiii.
						\item \dimd.
						\item \dimdd.
						\item \dimddd.
						\item \dimdddd
					\end{itemize}
				\end{minipage}
			}                                                            & \multirow{22}*[0cm]{
				\begin{minipage}[t]{6cm}
					\begin{itemize}[itemsep=-0pt,leftmargin=*,labelsep=.02cm]
						\item \textbf{Enfoque de estudio:} Investigación \MakeTextLowercase{\enfoque}.
						\item \textbf{Tipo de estudio:} Investigación \MakeTextLowercase{\tipo}. \item \textbf{Nivel de estudio:} Investigación \MakeTextLowercase{\nivel}.
						\item \textbf{Diseño de estudio:} Diseño \MakeTextLowercase{\diseno} de un mismo grupo con pre y postprueba.
						\item \textbf{Metodología:} Experimental, hipotético deductivo, estadístico y analítico.
						\item \textbf{Población:} Constituida por \poblacion. \item \textbf{Muestra:} Constituida por \muestra.
						\item \textbf{Muestreo:} Probabilístico. \item \textbf{Técnica:} Observación, evaluación peda\-gógica y experimentación.
						\item \textbf{Instrumentos:} Ficha de observación, ficha de opinión prueba escrita y módulos experimentales.
						\item \textbf{Procesamiento de datos:} Mediante los programas estadísticos R Sweave y Excel.
					\end{itemize}
				\end{minipage}
			} \Tstrut                                                                                                                                                                 \\ \cline{1-3} \problema&\objetivo&\hipotesis&&&\\ \cline{1-3} \ce\bf PROBLEMAS ESPECÍFICOS &\ce\bf OBJETIVOS ESPECÍFICOS &\ce\bf  HIPÓTESIS ESPECÍFICAS &&&\Tstrut\\\cline{1-3} \problemae&\objetivoe&\hipotesise&&&\\\cline{1-3}
			\problemaee                                                  & \objetivoee          & \hipotesisee     &                  &                      &                        \\\cline{1-3}
			\problemaeee                                                 & \objetivoeee         & \hipotesiseee    &                  &                      &                        \\\cline{1-3}
			\problemaeeee                                                & \objetivoeeee        & \hipotesiseeee   &                  &                      &                        \\ \hline
		\end{tabular}
	\end{table}
\end{landscape}



\begin{landscape}
	\subsection{Matriz de instrumentos}

	\lugar -- TÍTULO: \titulo -- AUTOR: \autor\\
	\begin{table}[ht!]\caption{Matriz de instrumentos de la variable \MakeTextLowercase{\variablei}}
		\centering  \renewcommand\tabcolsep{0.1cm}\renewcommand\arraystretch{1}
		\scriptsize
		\begin{tabular}{|c|l|l|l|c|c|}
			\hline
			\ce\bf VARIABLE                                        & \ce\bf DIMENSIONES        & \ce\bf INDICADORES                  & \ce\bf ÍTEMS & \ce\bf VALORACIÓN                                                                   & \bf INSTRUMENTOS \Tstrut                                         \\\hline
			\multirow{26}{*}{\rotatebox[origin=c]{90}{\variablei}} & \multirow{6}{*}{\dimi}    & \multirow{2}{*}{\gb}                &              & \multirow{26}{*}{\rotatebox[origin=c]{90}{Excelente Bueno Regular Malo Deficiente}} & \multirow{26}{*}{\rotatebox[origin=c]{90}{Ficha de observación}} \\\cline{4-4}
			                                                       &                           &                                     &              &                                                                                     &                                                                  \\\cline{3-4}
			                                                       &                           & \multirow{2}{*}{\gbb}               &              &                                                                                     &                                                                  \\\cline{4-4}
			                                                       &                           &                                     &              &                                                                                     &                                                                  \\\cline{3-4}
			                                                       &                           & \multirow{1}{*}{\gbbb}              &              &                                                                                     &                                                                  \\\cline{3-4}
			                                                       &                           & \multirow{1}{*}{\gbbbb}             &              &                                                                                     &                                                                  \\\cline{2-4}
			                                                       & \multirow{6}{*}{\dimii}   & \multirow{2}{*}{\gbbbbb}            &              &                                                                                     &                                                                  \\\cline{4-4}
			                                                       &                           &                                     &              &                                                                                     &                                                                  \\\cline{3-4}
			                                                       &                           & \multirow{1}{*}{\gbbbbbb}           &              &                                                                                     &                                                                  \\\cline{3-4}
			                                                       &                           & \multirow{1}{*}{\gbbbbbbb}          &              &                                                                                     &                                                                  \\\cline{3-4}
			                                                       &                           & \multirow{2}{*}{\gbbbbbbbb}         &              &                                                                                     &                                                                  \\\cline{4-4}
			                                                       &                           &                                     &              &                                                                                     &                                                                  \\\cline{2-4}
			                                                       & \multirow{7}{*}{\dimiii}  & \multirow{2}{*}{\gbbbbbbbbb}        &              &                                                                                     &                                                                  \\\cline{4-4}
			                                                       &                           &                                     &              &                                                                                     &                                                                  \\\cline{3-4}
			                                                       &                           & \multirow{1}{*}{\gbbbbbbbbbb}       &              &                                                                                     &                                                                  \\\cline{3-4}
			                                                       &                           & \multirow{2}{*}{\gbbbbbbbbbbb}      &              &                                                                                     &                                                                  \\\cline{4-4}
			                                                       &                           &                                     &              &                                                                                     &                                                                  \\\cline{3-4}
			                                                       &                           & \multirow{2}{*}{\gbbbbbbbbbbbb}     &              &                                                                                     &                                                                  \\\cline{4-4}
			                                                       &                           &                                     &              &                                                                                     &                                                                  \\\cline{2-4}
			                                                       & \multirow{7}{*}{\dimiiii} & \multirow{2}{*}{\gbbbbbbbbbbbbb}    &              &                                                                                     &                                                                  \\\cline{4-4}
			                                                       &                           &                                     &              &                                                                                     &                                                                  \\\cline{3-4}
			                                                       &                           & \multirow{1}{*}{\gbbbbbbbbbbbbbb}   &              &                                                                                     &                                                                  \\\cline{3-4}
			                                                       &                           & \multirow{2}{*}{\gbbbbbbbbbbbbbbb}  &              &                                                                                     &                                                                  \\\cline{4-4}
			                                                       &                           &                                     &              &                                                                                     &                                                                  \\\cline{3-4}
			                                                       &                           & \multirow{2}{*}{\gbbbbbbbbbbbbbbbb} &              &                                                                                     &                                                                  \\\cline{4-4}
			                                                       &                           &                                     &              &                                                                                     &                                                                  \\\hline
		\end{tabular}
	\end{table}
\end{landscape}

\begin{landscape}
	%\subsection{Matriz de instrumentos de la variable \MakeTextLowercase{\variabled} \label{1www}}

	TÍTULO: \titulo -- AUTOR: \autor
	\begin{table}[ht!]\caption{Matriz de instrumentos de la variable \MakeTextLowercase{\variabled}}
		\centering \scriptsize\renewcommand\tabcolsep{0.1cm}\renewcommand\arraystretch{1}
		\begin{tabular}{|c|l|l|l|c|c|}
			\hline
			\ce\bf VARIABLE                                        & \ce\bf DIMENSIONES        & \ce\bf INDICADORES                 & \ce\bf ÍTEMS                                                    & \ce\bf VALORACIÓN                                                                   & \bf INSTRUMENTOS \Tstrut                                         \\\hline
			\multirow{26}{*}{\rotatebox[origin=c]{90}{\variabled}} & \multirow{6}{*}{\dimd}    & \multirow{2}{*}{\fb}               & P1: ¿Utiliza variedad de colores?                               & \multirow{26}{*}{\rotatebox[origin=c]{90}{Excelente Bueno Regular Malo Deficiente}} & \multirow{26}{*}{\rotatebox[origin=c]{90}{Ficha de observación}} \\\cline{4-4}
			                                                       &                           &                                    & P2: ¿Plasma degradación entre los colores?                      &                                                                                     &                                                                  \\\cline{3-4}
			                                                       &                           & \multirow{2}{*}{\fbb}              & P3: ¿Utiliza formas y superficies regulares?                    &                                                                                     &                                                                  \\\cline{4-4}
			                                                       &                           &                                    & P4: ¿Utiliza formas irregulares someramente?                    &                                                                                     &                                                                  \\\cline{3-4}
			                                                       &                           & \multirow{1}{*}{\fbb}              & P3: ¿Utiliza formas y superficies regulares?                    &                                                                                     &                                                                  \\\cline{3-4}
			                                                       &                           & \multirow{1}{*}{\fbbb}             & P5: ¿Manifiesta degradación de texturas?                        &                                                                                     &                                                                  \\\cline{2-4}
			                                                       & \multirow{6}{*}{\dimdd}   & \multirow{2}{*}{\fbbbb}            & P6: ¿Distribuye adecuadamente los tonos?                        &                                                                                     &                                                                  \\\cline{4-4}
			                                                       &                           &                                    & P7: ¿Asocia adecuadamente los tonos?                            &                                                                                     &                                                                  \\\cline{3-4}
			                                                       &                           & \multirow{1}{*}{\fbbbbb}           & P8: ¿Asocia proporcionalmente los volúmenes?                    &                                                                                     &                                                                  \\\cline{3-4}
			                                                       &                           & \multirow{1}{*}{\fbbbbbb}          & P9: ¿Conjuga texturas diversas?                                 &                                                                                     &                                                                  \\\cline{3-4}
			                                                       &                           & \multirow{2}{*}{\fbbbbbbb}         & P9: ¿Conjuga texturas diversas?                                 &                                                                                     &                                                                  \\\cline{4-4}
			                                                       &                           &                                    & P10: ¿Homogeneiza texturas secuencialmente?                     &                                                                                     &                                                                  \\\cline{2-4}
			                                                       & \multirow{7}{*}{\dimddd}  & \multirow{2}{*}{\fbbbbbbbb}        & P11: ¿Representa elementos con dimensión creciente?             &                                                                                     &                                                                  \\\cline{4-4}
			                                                       &                           &                                    & P12: ¿Manifiesta elementos plásticos proporcionalmente?         &                                                                                     &                                                                  \\\cline{3-4}
			                                                       &                           & \multirow{1}{*}{\fbbbbbbbbb}       & P13: ¿Representa los elementos plásticos de manera decreciente? &                                                                                     &                                                                  \\\cline{3-4}
			                                                       &                           & \multirow{2}{*}{\fbbbbbbbbbb}      & P14: ¿Conjuga los elementos plásticos?                          &                                                                                     &                                                                  \\\cline{4-4}
			                                                       &                           &                                    & p15: ¿Distribuye elementos, constantes ligeramente?             &                                                                                     &                                                                  \\\cline{3-4}
			                                                       &                           & \multirow{2}{*}{\fbbbbbbbbbbb}     & P16: ¿Yuxtapone elementos plásticos mixtos?                     &                                                                                     &                                                                  \\\cline{4-4}
			                                                       &                           &                                    & P16: ¿Yuxtapone elementos plásticos mixtos?                     &                                                                                     &                                                                  \\\cline{2-4}
			                                                       & \multirow{7}{*}{\dimdddd} & \multirow{2}{*}{\fbbbbbbbbbbbb}    & P11: ¿Representa elementos con dimensión creciente?             &                                                                                     &                                                                  \\\cline{4-4}
			                                                       &                           &                                    & P12: ¿Manifiesta elementos plásticos proporcionalmente?         &                                                                                     &                                                                  \\\cline{3-4}
			                                                       &                           & \multirow{1}{*}{\fbbbbbbbbbbbbb}   & P13: ¿Representa los elementos plásticos de manera decreciente? &                                                                                     &                                                                  \\\cline{3-4}
			                                                       &                           & \multirow{2}{*}{\fbbbbbbbbbbbbbb}  & P14: ¿Conjuga los elementos plásticos?                          &                                                                                     &                                                                  \\\cline{4-4}
			                                                       &                           &                                    & p15: ¿Distribuye elementos, constantes ligeramente?             &                                                                                     &                                                                  \\\cline{3-4}
			                                                       &                           & \multirow{2}{*}{\fbbbbbbbbbbbbbbb} & P16: ¿Yuxtapone elementos plásticos mixtos?                     &                                                                                     &                                                                  \\\cline{4-4}
			                                                       &                           &                                    & P17: ¿Yuxtapone los ritmos en el conjunto?                      &                                                                                     &                                                                  \\\hline
		\end{tabular}
	\end{table}
\end{landscape}



\begin{landscape}
	\subsection{Ficha de observación de las variables}

	TÍTULO: \titulo -- AUTOR: \autor
	\begin{table}[ht!]\caption{Ficha de observación de la variable\MakeTextLowercase{\variablei}}
		\centering  \renewcommand\tabcolsep{0.1cm}\renewcommand\arraystretch{1}
		\scriptsize
		\begin{tabular}{|l|l|l|l|l|l|l|l|l|l|l|l|l|l|l|l|l|l|l|l|l|l|l|l|l|l|}
			\hline
			\multirow{2}{*}{ \ce\bf ÍTEMS} & \multicolumn{21}{c|}{ESTUDIANTES}\Tstrut                                                                                                                       \\\cline{2-22}
			                               & E1                                       & E2 & E3 & E4 & E5 & E6 & E7 & E8 & E9 & E10 & E11 & E12 & E13 & E14 & E15 & E16 & E17 & E18 & E19 & E20 & E2\Tstrut \\\hline
			                               &                                          &    &    &    &    &    &    &    &    &     &     &     &     &     &     &     &     &     &     &     &           \\\hline
			                               &                                          &    &    &    &    &    &    &    &    &     &     &     &     &     &     &     &     &     &     &     &           \\\hline
			                               &                                          &    &    &    &    &    &    &    &    &     &     &     &     &     &     &     &     &     &     &     &           \\\hline
			                               &                                          &    &    &    &    &    &    &    &    &     &     &     &     &     &     &     &     &     &     &     &           \\\hline
			                               &                                          &    &    &    &    &    &    &    &    &     &     &     &     &     &     &     &     &     &     &     &           \\\hline
			                               &                                          &    &    &    &    &    &    &    &    &     &     &     &     &     &     &     &     &     &     &     &           \\\hline
			                               &                                          &    &    &    &    &    &    &    &    &     &     &     &     &     &     &     &     &     &     &     &           \\\hline
			                               &                                          &    &    &    &    &    &    &    &    &     &     &     &     &     &     &     &     &     &     &     &           \\\hline
			                               &                                          &    &    &    &    &    &    &    &    &     &     &     &     &     &     &     &     &     &     &     &           \\\hline
			                               &                                          &    &    &    &    &    &    &    &    &     &     &     &     &     &     &     &     &     &     &     &           \\\hline
			                               &                                          &    &    &    &    &    &    &    &    &     &     &     &     &     &     &     &     &     &     &     &           \\\hline
			                               &                                          &    &    &    &    &    &    &    &    &     &     &     &     &     &     &     &     &     &     &     &           \\\hline
			                               &                                          &    &    &    &    &    &    &    &    &     &     &     &     &     &     &     &     &     &     &     &           \\\hline
			                               &                                          &    &    &    &    &    &    &    &    &     &     &     &     &     &     &     &     &     &     &     &           \\\hline
			                               &                                          &    &    &    &    &    &    &    &    &     &     &     &     &     &     &     &     &     &     &     &           \\\hline
			                               &                                          &    &    &    &    &    &    &    &    &     &     &     &     &     &     &     &     &     &     &     &           \\\hline
			                               &                                          &    &    &    &    &    &    &    &    &     &     &     &     &     &     &     &     &     &     &     &           \\\hline
			                               &                                          &    &    &    &    &    &    &    &    &     &     &     &     &     &     &     &     &     &     &     &           \\\hline
			                               &                                          &    &    &    &    &    &    &    &    &     &     &     &     &     &     &     &     &     &     &     &           \\\hline
			                               &                                          &    &    &    &    &    &    &    &    &     &     &     &     &     &     &     &     &     &     &     &           \\\hline
			                               &                                          &    &    &    &    &    &    &    &    &     &     &     &     &     &     &     &     &     &     &     &           \\\hline
			                               &                                          &    &    &    &    &    &    &    &    &     &     &     &     &     &     &     &     &     &     &     &           \\\hline
			                               &                                          &    &    &    &    &    &    &    &    &     &     &     &     &     &     &     &     &     &     &     &           \\\hline
			                               &                                          &    &    &    &    &    &    &    &    &     &     &     &     &     &     &     &     &     &     &     &           \\\hline
			                               &                                          &    &    &    &    &    &    &    &    &     &     &     &     &     &     &     &     &     &     &     &           \\\hline
			                               &                                          &    &    &    &    &    &    &    &    &     &     &     &     &     &     &     &     &     &     &     &           \\\hline
		\end{tabular}
	\end{table}
\end{landscape}




\begin{landscape}
	%\subsection{Matriz de instrumentos de la variable \MakeTextLowercase{\variabled} \label{1www}}

	TÍTULO: \titulo -- AUTOR: \autor
	\begin{table}[ht!]\caption{Ficha de observación de la variable \MakeTextLowercase{\variabled}}
		\centering \scriptsize\renewcommand\tabcolsep{0.1cm}\renewcommand\arraystretch{1}
		\begin{tabular}{|l|l|l|l|l|l|l|l|l|l|l|l|l|l|l|l|l|l|l|l|l|l|l|l|l|l|}
			\hline
			\multirow{2}{*}{ \ce\bf ÍTEMS} & \multicolumn{21}{c|}{ESTUDIANTES}\Tstrut                                                                                                                       \\\cline{2-22}
			                               & E1                                       & E2 & E3 & E4 & E5 & E6 & E7 & E8 & E9 & E10 & E11 & E12 & E13 & E14 & E15 & E16 & E17 & E18 & E19 & E20 & E2\Tstrut \\\hline
			                               &                                          &    &    &    &    &    &    &    &    &     &     &     &     &     &     &     &     &     &     &     &           \\\hline
			                               &                                          &    &    &    &    &    &    &    &    &     &     &     &     &     &     &     &     &     &     &     &           \\\hline
			                               &                                          &    &    &    &    &    &    &    &    &     &     &     &     &     &     &     &     &     &     &     &           \\\hline
			                               &                                          &    &    &    &    &    &    &    &    &     &     &     &     &     &     &     &     &     &     &     &           \\\hline
			                               &                                          &    &    &    &    &    &    &    &    &     &     &     &     &     &     &     &     &     &     &     &           \\\hline
			                               &                                          &    &    &    &    &    &    &    &    &     &     &     &     &     &     &     &     &     &     &     &           \\\hline
			                               &                                          &    &    &    &    &    &    &    &    &     &     &     &     &     &     &     &     &     &     &     &           \\\hline
			                               &                                          &    &    &    &    &    &    &    &    &     &     &     &     &     &     &     &     &     &     &     &           \\\hline
			                               &                                          &    &    &    &    &    &    &    &    &     &     &     &     &     &     &     &     &     &     &     &           \\\hline
			                               &                                          &    &    &    &    &    &    &    &    &     &     &     &     &     &     &     &     &     &     &     &           \\\hline
			                               &                                          &    &    &    &    &    &    &    &    &     &     &     &     &     &     &     &     &     &     &     &           \\\hline
			                               &                                          &    &    &    &    &    &    &    &    &     &     &     &     &     &     &     &     &     &     &     &           \\\hline
			                               &                                          &    &    &    &    &    &    &    &    &     &     &     &     &     &     &     &     &     &     &     &           \\\hline
			                               &                                          &    &    &    &    &    &    &    &    &     &     &     &     &     &     &     &     &     &     &     &           \\\hline
			                               &                                          &    &    &    &    &    &    &    &    &     &     &     &     &     &     &     &     &     &     &     &           \\\hline
			                               &                                          &    &    &    &    &    &    &    &    &     &     &     &     &     &     &     &     &     &     &     &           \\\hline
			                               &                                          &    &    &    &    &    &    &    &    &     &     &     &     &     &     &     &     &     &     &     &           \\\hline
			                               &                                          &    &    &    &    &    &    &    &    &     &     &     &     &     &     &     &     &     &     &     &           \\\hline
			                               &                                          &    &    &    &    &    &    &    &    &     &     &     &     &     &     &     &     &     &     &     &           \\\hline
			                               &                                          &    &    &    &    &    &    &    &    &     &     &     &     &     &     &     &     &     &     &     &           \\\hline
			                               &                                          &    &    &    &    &    &    &    &    &     &     &     &     &     &     &     &     &     &     &     &           \\\hline
			                               &                                          &    &    &    &    &    &    &    &    &     &     &     &     &     &     &     &     &     &     &     &           \\\hline
			                               &                                          &    &    &    &    &    &    &    &    &     &     &     &     &     &     &     &     &     &     &     &           \\\hline
			                               &                                          &    &    &    &    &    &    &    &    &     &     &     &     &     &     &     &     &     &     &     &           \\\hline
			                               &                                          &    &    &    &    &    &    &    &    &     &     &     &     &     &     &     &     &     &     &     &           \\\hline
			                               &                                          &    &    &    &    &    &    &    &    &     &     &     &     &     &     &     &     &     &     &     &           \\\hline
		\end{tabular}
	\end{table}
\end{landscape}

\subsection{Ficha de validación de instrumentos}
\begin{enumerate}
	\item DATOS GENERALES
	      \begin{itemize}
		      \item Apellidos y nombres del experto: \dotfill
		      \item Cargo e institución donde labora: \dotfill
		      \item Nombre de los instrumentos: Ficha de observación de la  variable indendiente (INST1) y Ficha de observación de la  variable dendiente (INST2)
		      \item Titulo de la tesis: \titulo
		      \item Autor de los instrumentos: \autor
	      \end{itemize}

	\item CRITERIOS DE VALIDACIÓN (Acuerdo: A, Descacuerdo: D)

	      \begin{center}\small
		      \begin{tabular}{|c|p{6.8cm}|c|c|c|c|}
			      \hline
			      \multirow{2}*[-.01cm]{INDICADORES} & \multirow{2}*[-.01cm]{CRITERIOS DE VALIDACIÓN}                                       & \multicolumn{2}{c|}{INST1} & \multicolumn{2}{c|}{INST2}         \\			      \cline{3-6}
			                                         &                                                                                      & A                          & D                          & A & D \\			      \hline
			      Claridad                           & ¿Está formulado con lenguaje claro, apropiado y cencillo?                            &                            &                            &   &   \\			      \hline
			      Objetividad                        & ¿Las preguntas realmente recojen datos de las variables y los indicadores?           &                            &                            &   &   \\			      \hline
			      Actualización                      & ¿El instrumento es adecuado para el tipo de variable en estudio?                     &                            &                            &   &   \\			      \hline
			      Organización                       & ¿La presentacion formal (tipo, tamaño, letra, etc.) del instrumento es apropiado?    &                            &                            &   &   \\			      \hline
			      Suficiencia                        & ¿Las preguntas son suficientes para recoger datos de todos los indicadores?          &                            &                            &   &   \\			      \hline
			      Intencionalidad                    & ¿Los items o preguntas responden al problema y objetivos del investigación?          &                            &                            &   &   \\			      \hline
			      Consistencia                       & ¿Los items o preguntas tienen sustento teórico y científico?                         &                            &                            &   &   \\			      \hline
			      Coherencia                         & ¿Los items o preguntas tienen son comprensibles y están bien redactados?             &                            &                            &   &   \\			      \hline
			      Metodología                        & ¿La estructura ofrece un orden logico y coherente por cada variable e indicador?     &                            &                            &   &   \\
			      \hline
			      Pertinencia                        & ¿El tipo de instrumento es pertinente para recoger datos de las variablesen estudio? &                            &                            &   &   \\
			      \hline
		      \end{tabular}
	      \end{center}

	\item OPINION DE APLICABILIDAD

	      ----------------------------------------------------------------------------------------------------------------------------------------------------------------
\end{enumerate}

\begin{center}

	---------------------------\\
	Firma del experto\\
	Telefono: .....................\\
	Fecha: .........................................
\end{center}


\end{document}